\documentclass[11pt]{amsart}
\usepackage[margin=1.15in]{geometry}
\usepackage[T1]{fontenc}
\usepackage{amsmath,amssymb,amsthm,mathtools,booktabs}
\usepackage[hidelinks]{hyperref}
\numberwithin{equation}{section}
\emergencystretch=3em

\newtheorem{theorem}{Theorem}[section]
\newtheorem{lemma}[theorem]{Lemma}
\newtheorem{proposition}[theorem]{Proposition}
\newtheorem{corollary}[theorem]{Corollary}
\theoremstyle{definition}
\newtheorem{definition}[theorem]{Definition}
\theoremstyle{remark}
\newtheorem{remark}[theorem]{Remark}

\newcommand{\li}{\operatorname{li}}
\newcommand{\Rp}{\operatorname{Re}}
\newcommand{\Ip}{\operatorname{Im}}
\newcommand{\eps}{\varepsilon}
\newcommand{\code}[1]{\texttt{#1}}
\DeclareMathOperator{\sgn}{sgn}

\title{Ramanujan's prime counting inequality holds for all $x \geq 38\,358\,837\,683$}
\author{Jaideep Sai Padhi}
\address{Purdue University, West Lafayette, IN 47907, USA}
\email{jpadhi@purdue.edu}
\subjclass[2020]{Primary 11N05; Secondary 11M26, 11Y35}
\keywords{Ramanujan's inequality, prime counting function, zero-free region, zero-density estimate, explicit formula}

\begin{document}

\begin{abstract}
Let $N_R$ denote the smallest positive integer such that Ramanujan's inequality
$\pi(x)^2 < \frac{ex}{\log x}\,\pi(x/e)$ holds for every real $x \geq N_R$.
Assuming the Riemann hypothesis, Dudek and Platt showed that $N_R = 38\,358\,837\,683$.
In this paper we prove the same value unconditionally.
For this purpose we close the range $e^{103} < x < e^{3158.442}$, which was not covered by previous unconditional results, using an explicit formula for $\psi(x)$ together with new explicit information on zeros of the Riemann zeta function slightly above the height $3 \cdot 10^{12}$, up to which the Riemann hypothesis has been verified.
\end{abstract}

\maketitle

\section{Introduction}\label{sec:intro}

Throughout, $\pi(x)$ is the prime counting function and $\li(x) = \int_0^x dt/\log t$ is the logarithmic integral (principal value).
The prime number theorem in the form $\pi(x) - \li(x) = O(x\exp(-c\sqrt{\log x}))$, with some $c > 0$, combined with the expansion $\li(x) = \frac{x}{\log x}\sum_{k<n}k!\log^{-k}x + O(x\log^{-n-1}x)$, which follows by integrating by parts, determines $\pi(x)$ up to an error that is smaller than any fixed power of $\log x$ times $x/\log x$.

Ramanujan noticed that this expansion forces the inequality
\begin{equation}\label{eq:ram}
\pi(x)^2 < \frac{ex}{\log x}\,\pi\Big(\frac{x}{e}\Big)
\end{equation}
for all sufficiently large $x$; see Berndt's edition of the notebooks \cite[p.~112]{BerndtIV}.
Indeed, inserting the expansion of $\li(x)$ into both sides of \eqref{eq:ram}, the terms up to order $x^2/\log^5 x$ cancel, and the right-hand side exceeds the left-hand side by $x^2/\log^6 x\,(1+o(1))$.
Hence \eqref{eq:ram} holds as soon as the relative error $|\pi(x) - \li(x)|\log x/x$ is smaller than about $1/(3\log^4 x)$.
As in \cite{Axler}, we write $N_R$ for the least positive integer with the property that \eqref{eq:ram} is true for all real $x \geq N_R$.

Under the assumption that the Riemann hypothesis is true, Hassani \cite{Hassani} gave an explicit value of $x_0$ such that \eqref{eq:ram} holds for every $x \geq x_0$.
Dudek and Platt \cite{DudekPlatt} combined a refined version of this argument with a direct computation for $x \leq 1.15\cdot 10^{16}$, and found that the largest integer counterexample to \eqref{eq:ram} is $x = 38\,358\,837\,682$; the largest prime counterexample, $38\,358\,837\,677$, had been found earlier by Galway (see \cite{DudekPlatt}).
Thus the Riemann hypothesis implies $N_R = 38\,358\,837\,683$.

Unconditionally, two ranges of $x$ have been settled, and the history of each is short.
For small $x$, Axler \cite{Axler2} showed, using B\"uthe's explicit bounds \cite{Buthe2018}, that \eqref{eq:ram} holds for $38\,358\,837\,683 \leq x \leq 10^{19}$; this range was extended to $x \leq e^{58}$ by Platt and Trudgian \cite{PlattTrudgian} and to $x \leq e^{103}$ by Johnston \cite[Theorem~5.1]{Johnston}.
For large $x$, the admissible lower limit $e^{X_0}$ has decreased from $X_0 = 9658$ (Dudek and Platt \cite{DudekPlatt}) through $9032$ \cite{Axler2}, $3915$ \cite{PlattTrudgian}, $3604$ \cite{CullyHugillJohnston} and $3361$ \cite{JohnstonYang} to $X_0 = 3158.442$ \cite[Theorem~1.1]{Axler}.
Consequently, unconditionally,
\[
38\,358\,837\,683 \leq N_R \leq e^{3158.442}.
\]
Two recent preprints of De \cite{De1,De2} claim to cover the remaining range; in Section~\ref{sec:remarks} we explain why the arguments given there do not prove \eqref{eq:ram}.

The main result of this paper closes the gap.

\begin{theorem}\label{thm:main}
The inequality \eqref{eq:ram} holds for every real $x$ with $e^{58} \leq x \leq e^{3159}$.
Consequently, $N_R = 38\,358\,837\,683$.
\end{theorem}

The second statement follows from the first one together with the results of Johnston \cite[Theorem~5.1]{Johnston} and Axler \cite[Theorem~1.1]{Axler} quoted above, and the computations of Galway and of Dudek and Platt \cite{DudekPlatt} for small $x$.

The only information about the zeros of the Riemann zeta function $\zeta(s)$ that we take from a large-scale computation is the theorem of Platt and Trudgian \cite{PlattTrudgianRH} that every non-trivial zero $\rho$ of $\zeta(s)$ with $0 < \Ip \rho \leq H_0$, where
\[
H_0 := 3\cdot 10^{12},
\]
satisfies $\Rp\rho = 1/2$.
Let us describe the method.
Write $L = \log x$.
By an elementary argument (Section~\ref{sec:prelim}), the inequality \eqref{eq:ram} at $x$ follows from bounds of the form $|\pi(y) - \li(y)| \leq \eps_R(L)\,y/\log y$ at $y = x$ and $y = x/e$, where $\eps_R(L)$ is an explicit function with $\eps_R(L) \sim 1/(3L^4)$, and these bounds in turn follow from relative bounds for $\psi(t) - t$.
To bound $\psi(t) - t$ we use a smoothed explicit formula with Logan's kernel, following B\"uthe \cite{Buthe2016,Buthe2018} (Section~\ref{sec:explicit}).
The zeros of height at most $H_0$ lie on the critical line, while the zeros above $H_0$ are controlled by a zero-free region and by zero-density estimates.

The binding range turns out to be $L \approx 900$.
There the kernel has bandwidth about $e^{30}$, and the decisive zeros are hypothetical zeros $\beta + i\gamma$ with $(1-\beta)\log\gamma \approx 0.21$ at heights just above $H_0$.
With the published zero-free regions and zero-density estimates alone, the contribution of these zeros exceeds the available error budget several times over.
The present paper therefore supplies three new explicit inputs.
\begin{enumerate}
\item In Section~\ref{sec:zfr} we evaluate the inequality behind the zero-free region of Bellotti, Trudgian and Yang \cite{BTY} exactly at every height, instead of through a single uniform constant.
This gives $(1-\beta)\log\gamma \geq 0.21384$ for zeros with $H_0 < \gamma \leq e^{29}$, that is, the constant $4.676$ in place of $4.896$ in the zero-free region $\sigma > 1 - 1/(R\log t)$.
\item In Section~\ref{sec:rep} we keep a second zero in the same inequality and show that two zeros close to the line $\sigma = 1$ cannot be close to each other: for example, two zeros $\beta_j + i\gamma_j$ with $H_0 < \gamma_j \leq e^{29}$ and $(1-\beta_j)\log\gamma_j \leq 0.348$ satisfy $|\gamma_0 - \gamma_1| \geq 0.058$.
\item In Section~\ref{sec:density} we rework Bellotti's log-free zero-density estimate \cite{Bellotti} for $\sigma$ close to $1$.
The main gain comes from an explicit and asymptotically sharp evaluation of the Barban--Vehov mean square, which is the quantity Bellotti denotes by $d_4$.
Our bound for $d_4$ is about $1.5$ at our parameters, while her Table~5 gives values between $85$ and $105$ at hers.
\end{enumerate}
The proof of Theorem~\ref{thm:main} is completed in Section~\ref{sec:proof} by a computation, which is carried out entirely in rigorous interval (ball) arithmetic using the Arb library \cite{arb}.
Floating-point arithmetic is used only to choose parameters, and every choice yields a valid bound.
The code and the output of all computations are available at \url{https://github.com/jaideepsaipadhi/ramanujans-inequality}.
While re-deriving the explicit inputs, we found a few inaccuracies in the literature; these are collected in Section~\ref{sec:remarks}.
None of them affects the main results of the papers concerned.

\section{Preliminaries}\label{sec:prelim}

In this section we reduce Theorem~\ref{thm:main} to explicit bounds for $\psi(t) - t$.
For $L > 1$ and $x = e^L$ we put
\begin{equation}\label{eq:gb}
g(L) = \frac{\li(x)^2 - \frac{ex}{L}\li(x/e)}{(x/L)^2}, \qquad b(L) = \frac{2\li(x)}{x/L} + \frac{L}{L-1},
\end{equation}
and, whenever $g(L) < 0$,
\begin{equation}\label{eq:epsR}
\eps_R(L) = \frac{-b(L) + \sqrt{b(L)^2 - 4g(L)}}{2}.
\end{equation}
Then $\eps_R(L)$ is the positive root of $\eps^2 + b(L)\eps + g(L) = 0$.
Asymptotically, $\eps_R(L) = \frac{1}{3L^4}(1 + O(L^{-1}))$; for instance, $3L^4\eps_R(L)$ equals $1.2706$ at $L = 58$, $1.0145$ at $L = 907$ and $1.0041$ at $L = 3159$.

\begin{lemma}\label{lem:reduce}
Let $L = \log x \geq 58$ with $g(L) < 0$, and let $\eps = \eps_R(L)$.
Assume that
\begin{equation}\label{eq:pibounds}
|\pi(x) - \li(x)| \leq \eps\,\frac{x}{L} \quad\text{and}\quad \Big|\pi\Big(\frac xe\Big) - \li\Big(\frac xe\Big)\Big| \leq \eps\,\frac{x/e}{L-1},
\end{equation}
where at least one of the two inequalities is strict.
Then \eqref{eq:ram} holds at $x$.
\end{lemma}

\begin{proof}
Since $\li(x) > \eps x/L$, the inequalities \eqref{eq:pibounds} give $0 < \pi(x) \leq \li(x) + \eps x/L$ and $\pi(x/e) \geq \li(x/e) - \eps x/(e(L-1))$.
Hence it suffices to show that
\[
\Big(\li(x) + \eps\frac{x}{L}\Big)^2 < \frac{ex}{L}\Big(\li\Big(\frac xe\Big) - \eps\frac{x}{e(L-1)}\Big).
\]
After division by $(x/L)^2$, this inequality reads $g(L) + b(L)\eps + \eps^2 < 0$.
For $\eps = \eps_R(L)$ the left-hand side vanishes, and the strict inequality in \eqref{eq:pibounds} makes the inequality strict.
\end{proof}

The direct evaluation of $\eps_R(L)$ requires some care, since $g(L)$ is the difference of two quantities of size $1$ which agree up to $O(L^{-4})$.
The following lemma provides an evaluation without cancellation.

\begin{lemma}\label{lem:epsR}
Let $A(L) = Le^{-L}\li(e^L)$, $u = 1/L$ and $P(u) = \sum_{k=0}^{15}k!\,u^k$.
Then
\[
g(L) = A(L)^2 - \frac{L}{L-1}A(L-1), \qquad b(L) = 2A(L) + \frac{L}{L-1},
\]
and $\eps_R(L) = -2g(L)/\big(b(L) + \sqrt{b(L)^2 - 4g(L)}\big)$.
Moreover, $A(L) = P(u) + R(L)$, where
\begin{equation}\label{eq:Rbound}
-6e^{-L/2} \leq R(L) \leq \frac{16!\,u^{16}}{1 - 34u} + 6e^{-L/2} \qquad (L \geq 57).
\end{equation}
In particular, $g(L) < 0$ for every $L \in [58, 3160.1]$, and on every interval $[a, a + 0.05] \subseteq [58, 3160.1]$ the function $\eps_R$ has an explicit lower bound which we compute in ball arithmetic.
\end{lemma}

\begin{proof}
The formulas for $g$, $b$ and $\eps_R$ follow from \eqref{eq:gb} and \eqref{eq:epsR}.
To prove \eqref{eq:Rbound}, we integrate $\int_a^x dt/\log t$ by parts $16$ times, where $a = \sqrt x$.
Since $A$ and $P$ are both bounded by $1.1$ for arguments at least $28.5$, the boundary term at $a$, namely $\li(a) - aP(2/L)/(L/2)$, is at most $3a/(L/2)$ in absolute value.
The logarithmic derivative of $e^s s^{-17}$ is at least $1 - 34/L$ on $[L/2, L]$, so the remaining integral satisfies
\[
0 \leq \frac{L}{x}\,16!\int_a^x\frac{dt}{\log^{17}t} \leq \frac{16!\,u^{16}}{1 - 34u}.
\]
This proves \eqref{eq:Rbound}.
Next, $S(u) = P(u)^2 - P(u/(1-u))/(1-u)$ is a power series in $u$ with integer coefficients $s_n$, where $s_0 = s_1 = s_2 = s_3 = 0$, $s_4 = -1$ and $s_5 = -14$.
We compute these coefficients exactly up to $n = 70$; since the coefficient of $u^n$ in $P(u/(1-u))/(1-u)$ is $\sum_kk!\binom nk \leq 2^n\sum_{k\leq15}k!$, and similarly for $P(u)^2$, the tail is at most $(\sum_{k\leq15}k!)(2u)^{71}/(1-2u)$.
Writing $g(L)$ as $u^4\cdot(S(u)/u^4)$ plus terms involving only $R$, we evaluate it without cancellation in ball arithmetic on each interval $[a, a+0.05]$.
The enclosures for $\eps_R$ obtained in this way have relative width at most $1.9\cdot10^{-3}$.
\end{proof}

We write $y = x/e$.
Then $\log y = L - 1$, and Lemma~\ref{lem:reduce} shows that Theorem~\ref{thm:main} follows from the two inequalities
\begin{equation}\label{eq:target}
\begin{aligned}
|\pi(y) - \li(y)| &< \eps_R(\log y + 1)\,\frac{y}{\log y} \qquad (57 \leq \log y \leq 3158),\\
|\pi(y) - \li(y)| &< \eps_R(\log y)\,\frac{y}{\log y} \qquad\quad\ \, (58 \leq \log y \leq 3159).
\end{aligned}
\end{equation}
We do not use any monotonicity of $\eps_R$: both inequalities will be checked directly on small intervals.

To deduce \eqref{eq:target} from bounds for $\psi$, we use Chebyshev's function $\vartheta(x) = \sum_{p\leq x}\log p$ and the following identity.

\begin{lemma}\label{lem:transfer}
For every $y \geq 2$, we have
\begin{equation}\label{eq:pili}
\pi(y) - \li(y) = \frac{\vartheta(y) - y}{\log y} + \int_2^y \frac{\vartheta(t) - t}{t\log^2t}\,dt + \frac{2}{\log 2} - \li(2),
\end{equation}
where $|2/\log 2 - \li(2)| < 1.8403$.
Moreover, $0 \leq \psi(t) - \vartheta(t) \leq 1.4262\sqrt t$ for every $t > 0$.
\end{lemma}

\begin{proof}
By partial summation, $\pi(y) = \vartheta(y)/\log y + \int_2^y \vartheta(t)\,dt/(t\log^2t)$, and integration by parts gives $\li(y) = y/\log y + \int_2^y dt/\log^2t + \li(2) - 2/\log 2$.
Subtracting these identities yields \eqref{eq:pili}.
The bound for $\psi(t) - \vartheta(t)$ is due to Rosser and Schoenfeld \cite[Theorem~13]{RS62}.
\end{proof}

In Section~\ref{sec:proof} we cover $[43, 3159]$ by closed intervals $[L_1, L_2]$ and show, for each of them, an inequality of the form
\begin{equation}\label{eq:psicert}
|\psi(t) - t| \leq q_{[L_1,L_2]}\,\frac{\eps_R^-(L_2+1)}{1 + 2/L_1}\,t + a(t) \qquad (\log t \in [L_1, L_2]),
\end{equation}
where $\eps_R^-(L_2 + 1)$ is the lower endpoint of our enclosure of $\eps_R(L_2+1)$, $q_{[L_1,L_2]}$ is an explicit number, and
\begin{equation}\label{eq:at}
a(t) = 2 + e^{2\cdot10^{-4}}\big(\log t + 10^{-4}\big)\Big(2.01\cdot10^{-4}\sqrt t + \frac12\log\log(2t^2)\Big).
\end{equation}
The factor $(1 + 2/L_1)^{-1}$ leaves room for the integral in \eqref{eq:pili}.

\begin{proposition}\label{prop:transfer}
Let \eqref{eq:psicert} hold on every interval of the cover of $[43, 3159]$ constructed in Section~\ref{sec:proof}, with the values $q_{[L_1,L_2]}$ obtained there.
Then, for every $y$ with $57 \leq \log y \leq 3159$, the bound for $|\pi(y) - \li(y)|$ obtained from \eqref{eq:pili} is at most $0.99052\min\{\eps_R(\log y + 1), \eps_R(\log y)\}\,y/\log y$, where the second argument is used only for $\log y \geq 58$.
In particular, \eqref{eq:target} holds.
\end{proposition}

\begin{proof}
We use Lemma~\ref{lem:transfer}, the inequality $|\vartheta(t) - t| \leq |\psi(t) - t| + 1.4262\sqrt t$, and B\"uthe's bound $|\psi(t) - t| \leq 0.94\sqrt t$ for $11 < t \leq 10^{19}$ \cite[Theorem~2]{Buthe2018}.
For $t \leq 11$ we use the trivial bound $|\vartheta(t) - t| \leq 11$.
Hence the part of the integral in \eqref{eq:pili} over $[2, e^{43}]$ is at most
\[
\frac{2\,(0.94 + 1.4262)\,e^{21.5}}{\log^2 2} + 121.
\]
On $[e^{43}, y]$ we substitute $t = e^v$ and bound the integral by upper sums with step $0.05$ in $v$; on each step we use the largest bound \eqref{eq:psicert} that meets it, together with $\int_a^b e^v v^{-2}\,dv \leq (e^b - e^a)/a^2$.
For $y$ with $\log y \in [a, a + 0.05]$ and $a \geq 57$, the resulting bound for $|\pi(y) - \li(y)|\log y/y$ is compared with the lower bounds of Lemma~\ref{lem:epsR} for $\eps_R$ on $[a+1, a+1.05]$ and, if $a \geq 58$, on $[a, a+0.05]$.
This is carried out in ball arithmetic for every such interval $[a, a+0.05] \subseteq [57, 3159]$, and the largest ratio, $0.99052$, occurs at $\log y = 877$.
Note that $q_{[L_1,L_2]} < 1$ alone would not suffice, because of the additional terms $1.4262\sqrt t$, $a(t)$ and the integral in \eqref{eq:pili}; the proposition rests on this computation.
\end{proof}

Heuristically, the factor $(1 + 2/L_1)^{-1}$ in \eqref{eq:psicert} is explained as follows: since $\int^y\eps_R(\log t + 1)\,dt/\log^2t \approx \eps_R(\log y + 1)\,\frac{y}{\log^2y}(1 + \frac{6}{\log y})$, the integral in \eqref{eq:pili} adds a relative amount of about $1/L + 6/L^2 \leq 2/L$.
The proof of Proposition~\ref{prop:transfer} does not use this heuristic.

\section{An explicit formula for \texorpdfstring{$\psi(x)$}{psi(x)}}\label{sec:explicit}

We use the smoothing introduced by B\"uthe \cite{Buthe2016,Buthe2018}, but give our own proof of the smoothing step (Lemma~\ref{lem:smooth}); see Remark~\ref{rem:buthe} for the reason.
For $c > 0$ and $0 < \eps \leq 10^{-4}$ let
\[
\eta_{c,\eps}(\tau) = \frac{c}{\eps\sinh c}\,I_0\Big(c\sqrt{1 - (\tau/\eps)^2}\Big)\mathbf 1_{|\tau|\leq\eps}, \qquad
\ell_{c,\eps}(z) = \frac{c}{\sinh c}\cdot\frac{\sinh\sqrt{c^2 - (\eps z)^2}}{\sqrt{c^2 - (\eps z)^2}},
\]
where $I_0$ denotes the modified Bessel function of order $0$.
Then $\int\eta_{c,\eps} = 2$ and $\int\eta_{c,\eps}(\tau)e^{-iz\tau}\,d\tau = 2\ell_{c,\eps}(z)$.
Logan's function $\ell_{c,\eps}$ is entire and even with $\ell_{c,\eps}(0) = 1$.
Let $\lambda = \int\eta_{c,\eps}(\tau)e^{\tau/2}\,d\tau = 2\ell_{c,\eps}(i/2) \geq 2$ and
\begin{equation}\label{eq:phi}
\phi(\tau) = \frac{\eta_{c,\eps}(\tau)e^{\tau/2}}{\lambda}, \qquad \psi_{c,\eps}(x) = \int\psi_0(xe^{-\tau})\phi(\tau)\,d\tau,
\end{equation}
where $\psi_0$ is the normalised Chebyshev function.
So $\phi$ is a probability density on $[-\eps,\eps]$, and $\psi_{c,\eps}(x)$ is a weighted mean of $\psi_0$ over $[xe^{-\eps}, xe^{\eps}]$; it coincides with B\"uthe's function \cite[(3.2)--(3.4)]{Buthe2018}.

\begin{lemma}\label{lem:smooth}
Let $x \geq 10^{18}$, let $c \in \{30, 31, \ldots, 110\}$, let $0 < \eps \leq 10^{-4}$ and put $\lambda_s = \log(10\sqrt c/\eps)$.
Assume that $\lambda_s < \log x - 1$, and let $\nu_c = \int_0^1 u\,\eta_{c,1}(u)\,du$.
Then
\begin{equation}\label{eq:smooth}
|\psi(x) - \psi_{c,\eps}(x)| \leq \frac{\log x + \lambda_s}{\log x - \lambda_s}\cdot 2\eps e^{\eps/2}\Big(\frac{\nu_c}2 + \frac{I_0(c)}{200\sinh c}\Big)x + e^{2\eps}\log(e^\eps x)\Big(2.01\eps\sqrt x + \frac12\log\log(2x^2)\Big).
\end{equation}
Moreover, $2\sqrt c\,e^{\eps/2}(\nu_c/2 + I_0(c)/(200\sinh c)) \leq 0.80319 < 0.7946\cdot 1.02$ for all these $c$.
\end{lemma}

\begin{proof}
Since $\phi$ is a probability density and $\psi$ is non-decreasing, we have
\[
\psi(x) - \psi_{c,\eps}(x) \leq \int_0^\eps\big(\psi(x) - \psi((xe^{-\tau})^-)\big)\phi(\tau)\,d\tau,
\]
because the integrand is non-positive for $\tau < 0$.
For $3 \leq Y \leq X/2$, the Brun--Titchmarsh inequality $\pi(X) - \pi(X-Y) \leq 2Y/\log Y$ of Montgomery and Vaughan \cite{MV73}, together with the trivial count of prime powers as in \cite[(3.9)]{Buthe2018}, gives
\[
\psi(X) - \psi(X - Y) \leq \log X\Big(\frac{2Y}{\log Y} + 0.6\frac{Y}{\sqrt X} + \log\log X^2\Big).
\]
We apply this with $Y = x(1 - e^{-\tau}) \leq x\tau$.
For $\tau \geq \tau_0 := \eps/(10\sqrt c)$ we have $\log Y \geq \log x - \lambda_s$, and for $\tau < \tau_0$ we bound the integrand by its value at $\tau_0$.
The mass of $\phi$ on $[0, \tau_0]$ is at most $\tau_0\max\phi \leq \tau_0e^{\eps/2}cI_0(c)/(2\eps\sinh c)$, and $\int_0^\eps\tau\phi(\tau)\,d\tau \leq e^{\eps/2}\eps\nu_c/2$ since $\lambda \geq 2$.
This gives the upper bound for $\psi(x) - \psi_{c,\eps}(x)$.
The lower bound is obtained in the same way from $\psi_{c,\eps}(x) - \psi(x) \leq \int_{-\eps}^0(\psi_0(xe^{|\tau|}) - \psi(x))\phi(\tau)\,d\tau$; here the intervals end at $xe^{\eps}$, which leads to the factor $\log(e^\eps x) \leq \log x + \lambda_s$.
The prime powers and the term $\log\log X^2$ give the last term in \eqref{eq:smooth}.
Finally, the numerical inequality for the constant is verified in ball arithmetic for each $c$, with $\nu_c$ computed by rigorous numerical integration; the maximum occurs at $c = 110$.
\end{proof}

The next result is the explicit formula for $\psi_{c,\eps}$ together with a bound for the contribution of zeros of large height.

\begin{proposition}[{B\"uthe \cite[Propositions~2 and 3]{Buthe2016}, \cite[Proposition~2]{Buthe2018}}]\label{prop:explicit}
Let $x \geq 10$, $0 < \eps \leq 10^{-4}$ and $c \geq 3$.
Then
\[
x - \psi_{c,\eps}(x) = \frac{1}{\ell_{c,\eps}(i/2)}\sum_\rho\ell_{c,\eps}\Big(\frac{\rho - 1/2}{i}\Big)\frac{x^\rho}{\rho} + \Theta,
\]
where $|\Theta| \leq 2$, and
\[
\sum_{|\Ip\rho| > c/\eps}\Big|\frac{\ell_{c,\eps}((\rho - \frac12)/i)}{\ell_{c,\eps}(i/2)}\cdot\frac{x^\rho}{\rho}\Big| \leq 0.16\,\frac{x+1}{\sinh c}\,e^{0.71\sqrt{c\eps}}\log(3c)\log\frac c\eps.
\]
\end{proposition}

\begin{lemma}\label{lem:logancomplex}
Let $\rho = \beta + i\gamma$ with $0 < \beta < 1$ and $|\eps\gamma| \leq c$.
Then $|\ell_{c,\eps}((\rho - 1/2)/i)| \leq e^{\eps(c+1)/6}\ell_{c,\eps}(\gamma)$ and $\ell_{c,\eps}(i/2) \geq 1$.
\end{lemma}

\begin{proof}
We write $\ell_{c,\eps}(z) = \frac{c}{\sinh c}G(c^2 - \eps^2z^2)$ with $G(w) = \sum_{n\geq0}w^n/(2n+1)!$.
For $z = \gamma - i(\beta - \frac12)$ we get $|c^2 - \eps^2z^2| \leq c^2 - \eps^2\gamma^2 + \eps(c+1)$.
Since $G$ has positive coefficients, $|G(w)| \leq G(|w|)$, and since $\frac{d}{dw}\log G(w) \leq 1/6$ for $w \geq 0$, the first claim follows.
The second claim holds because $G$ is increasing on $[0,\infty)$.
\end{proof}

\begin{lemma}\label{lem:sumrecip}
For every $T \geq 10^4$, we have $\sum_{0 < \gamma \leq T}1/\gamma \leq \frac12\log^2T$, where the sum runs over the zeros of $\zeta(s)$ counted with multiplicity.
\end{lemma}

\begin{proof}
Let $N(T)$ denote the number of zeros with $0 < \gamma \leq T$.
We have $\sum_{0<\gamma\leq T}1/\gamma = N(T)/T + \int_{14}^T N(t)t^{-2}\,dt$.
Inserting the bound $|N(t) - \frac{t}{2\pi}\log\frac{t}{2\pi e} - \frac78| \leq 0.1038\log t + 0.2573\log\log t + 9.3675$ of Hasanalizade, Shen and Wong \cite{HSW}, we obtain $\sum_{0<\gamma\leq T}1/\gamma \leq \frac{1}{4\pi}\log^2T + \log T + 2$, which is smaller than $\frac12\log^2T$ for $T \geq 10^4$.
\end{proof}

From now on let $H = H_0 = 3\cdot10^{12}$ and $T_1 = c/\eps$, and for $s \in [0.6, 1)$ let
\[
S(s) = e^{\eps(c+1)/6}\sum_{\substack{H < \gamma \leq T_1\\ \beta > s}}\frac{\ell_{c,\eps}(\gamma)}{\gamma}.
\]
Further, let $N(\sigma, T)$ denote the number of zeros $\beta + i\gamma$ of $\zeta(s)$ with $\beta > \sigma$ and $0 < \gamma \leq T$, counted with multiplicity.
The following proposition is the basis of the computation in Section~\ref{sec:proof}.

\begin{proposition}\label{prop:master}
Let $43 \leq L_1 < L_2$, $c \in \{30, \ldots, 110\}$ and $0 < \eps \leq 10^{-4}$ with $\lambda_s = \log(10\sqrt c/\eps) < L_1 - 1$.
Assume that every zero $\beta + i\gamma$ with $H < \gamma \leq T_1$ satisfies $\beta \leq s_{\max}$.
Then, for every $x$ with $\log x \in [L_1, L_2]$,
\begin{multline}\label{eq:master}
\frac{|\psi(x) - x|}{x} \leq \frac{L_2 + \lambda_s}{L_1 - \lambda_s}\cdot\frac{0.7946\cdot1.02\,\eps}{\sqrt c} + \frac{0.32\,e^{0.71\sqrt{c\eps}}\log(3c)\log T_1}{\sinh c} + e^{-L_1/2}\log^2H\\
+ 2\Big(L_2\int_{0.6}^{s_{\max}}e^{(s-1)L_1}S(s)\,ds + e^{-0.4L_1}\big(S(0.6) + \log^2T_1\big)\Big) + \frac{a(x)}{x},
\end{multline}
where $a(x)$ is given by \eqref{eq:at}.
\end{proposition}

\begin{proof}
We combine Lemma~\ref{lem:smooth} with Proposition~\ref{prop:explicit}.
Since $0.16(x+1) \leq 0.32x$, the zeros with $|\gamma| > T_1$ contribute the second term.
The zeros with $|\gamma| \leq H$ lie on the critical line \cite{PlattTrudgianRH}, so by Lemmas~\ref{lem:logancomplex} and~\ref{lem:sumrecip} they contribute at most $2x^{1/2}\cdot\frac12\log^2H$.
The zeros with $H < |\gamma| \leq T_1$ occur in pairs $\rho, \bar\rho$, which gives the factor $2$.
For those with $\beta > 0.6$ we use the identity
\[
e^{\eps(c+1)/6}\sum_{\substack{H<\gamma\leq T_1\\ \beta > 0.6}}x^{\beta-1}\frac{\ell_{c,\eps}(\gamma)}{\gamma} = x^{-0.4}S(0.6) + \log x\int_{0.6}^1x^{s-1}S(s)\,ds,
\]
together with Lemma~\ref{lem:logancomplex}, the inequalities $x^{s-1} \leq e^{(s-1)L_1}$ and $\log x \leq L_2$, and $S(s) = 0$ for $s \geq s_{\max}$.
All remaining zeros in this range of heights have $\beta \leq 0.6$; they contribute at most $x^{-0.4}e^{\eps(c+1)/6}\sum_{H<\gamma\leq T_1}1/\gamma \leq x^{-0.4}\log^2T_1$, since Lemma~\ref{lem:sumrecip} gives $\frac12\log^2T_1$ and $e^{\eps(c+1)/6} \leq e^{0.002} < 2$.
The terms $\Theta$ and those of Lemma~\ref{lem:smooth} give $a(x)/x$.
\end{proof}

To bound $S(s)$ we use summation by parts.

\begin{lemma}\label{lem:stieltjes}
Let $f(t) = \ell_{c,\eps}(t)/t$ for $\eps t < c$ and $f(t) = c/(t\sinh c)$ for $\eps t \geq c$, and let $H = t_0 < t_1 < \cdots < t_n$ with $t_n \geq T_1$.
Let $N^+(s,t)$ be any upper bound for $N(s,t)$, and assume that $\lambda^*(l)/l$ is non-increasing, where every zero with $\gamma > H$ satisfies $(1-\beta)\log\gamma \geq \lambda^*(\log\gamma)$.
If we replace $N^+(s, t_{i+1})$ by $0$ whenever $s \geq 1 - \lambda^*(\log t_{i+1})/\log t_{i+1}$, then
\[
S(s)e^{-\eps(c+1)/6} \leq \sum_{i=0}^{n-1}N^+(s,t_{i+1})\big(f(t_i) - f(t_{i+1})\big) + N^+(s,t_n)f(t_n).
\]
\end{lemma}

\begin{proof}
The function $f$ is positive and non-increasing on $(0,\infty)$; enlarging the range of heights from $T_1$ to $t_n$ can only increase the sum.
Zeros with $\gamma \leq H$ do not satisfy $\beta > s \geq 0.6$.
Hence the inequality follows by summation by parts.
Since $\lambda^*(l)/l$ is non-increasing, $s \geq 1 - \lambda^*(\log t_{i+1})/\log t_{i+1}$ implies $N(s, t_{i+1}) = 0$.
\end{proof}

\section{A zero-free region depending on the height}\label{sec:zfr}

Bellotti, Trudgian and Yang \cite[Theorem~1]{BTY} proved that $\zeta(\sigma + it) \neq 0$ for $\sigma > 1 - 1/(4.896\log t)$ and $t \geq 3$.
The constant $4.896$ is determined by heights near $e^{76}$.
At the heights relevant to us, namely $\log t \leq 43$, their basic inequality permits a smaller constant.
In this section we evaluate it exactly at every height.
Since several constants in \cite{BTY} are computed only for $\sigma \geq 0.99346$, we first re-derive them for the smaller values of $\sigma$ and the larger values of $\eta$ that we need.

\subsection{The inequality of Bellotti, Trudgian and Yang}

We keep the notation of \cite{BTY}.
Let $\theta = 1.1338$, and let $w$ be the function of Heath-Brown type defined in \cite[Definition~1]{BTY}, which is supported on $[0,D]$ with $D = 2\theta\cot\theta$.
Let $W(z) = \int_0^De^{-zu}w(u)\,du$ and $w_0 = w(0) = 5.672787598\ldots$.
Let $\kappa_m$ $(0 \leq m \leq 6)$ be the coefficients in \cite[(18)]{BTY}, with $\kappa = \sum_m\kappa_m = 433/859$, and let $P(x) = \sum_{k\leq16}a_k\cos(kx) \geq 0$ be the trigonometric polynomial of \cite[Section~4]{BTY}, with $a_0 = 1$, $a_1 = 865534/497079$ and $a = \sum_{k\geq1}a_k = 2919857/828465$.
Finally, let $K = 16$ and $T_0 = 10^{10}$.
For $\eta > 0$ and $\frac12 < \sigma < 1$ we put
\[
f(u) = \eta w(\eta u)\sum_m\kappa_me^{-(2\sigma-1)mu}, \qquad F(z) = \sum_m\kappa_mW\Big(\frac{z + (2\sigma - 1)m}{\eta}\Big),
\]
so that $f(0) = \kappa\eta w_0$.

Let $\rho_0 = \beta_0 + it$ be a zero with $t \geq H_0$, and put $\eta = 1 - \beta_0$ and $\lambda_0 = \eta\log t$.
Suppose that $\zeta(s)$ has no zeros $\beta + i\gamma$ with $\beta \geq 1 - A/\log\gamma$ and $H_0 \leq \gamma \leq Kt + T_0$, and put $\sigma = 1 - A/\log(Kt + T_0)$.
Assume that $\eta \leq \eta_0 \leq 0.0145$ and $\sigma \geq \sigma_0 \geq 0.993$.
Then the argument of \cite[Lemma~15]{BTY} gives
\begin{multline}\label{eq:bty}
c_{10}\,\eta w_0 \leq \eta w_0\Big(\frac{a\kappa}2\log t - G\Big) + a_0F(\sigma - 1) - a_1\big(F(\sigma - 1 + \eta) + F(\sigma - \eta)\big) - Z\\
+ 10^{-10}\eta + 2\cdot10^{-7} + E_2\lambda_0\eta^2 + E_3\eta^3,
\end{multline}
where $Z \geq 0$ is the part of the sum over zeros that is discarded, and the constants $c_{10}$, $G$, $E_2$, $E_3$ depend on $(\eta_0, \sigma_0)$ and are given in Lemmas~\ref{lem:L89}--\ref{lem:L12} below.
Consequently, there is no such zero as soon as
\begin{multline}\label{eq:margin1}
\mathcal M_1 := a_1\big(F(\sigma - 1 + \eta) + F(\sigma - \eta)\big)\\
- \Big[\eta w_0\Big(\frac{a\kappa}{2}\log t - G - c_{10}\Big) + a_0F(\sigma - 1) + 10^{-10}\eta + 2\cdot10^{-7} + E_2\lambda_0\eta^2 + E_3\eta^3\Big]
\end{multline}
is positive.
In \cite{BTY}, the term coming from their Lemma~10 enters the final inequality with $f(0) = \eta w_0$ instead of $\kappa\eta w_0$; see Remark~\ref{rem:btyf0}.

\subsection{The constants}

\begin{lemma}\label{lem:W0}
Let $W_0(z) = W(z) - w_0/z$, and let $\nu, r > 0$.
For $\Rp z \geq \nu$ and $|z| \geq r$ we have $|W_0(z)| \leq C_w(\nu, r)|z|^{-3}$, where
\[
C_w(\nu, r) = |w''(0)| + \frac1r\Big(|w'''(0)| + e^{-\nu D}|w'''(D)| + \int_0^D|w''''(u)|e^{-\nu u}\,du\Big),
\]
$|w''(0)| = 50.06970149\ldots$ and $|w'''(0)| = |w'''(D)| = 117.26477652\ldots$.
\end{lemma}

\begin{proof}
We have $w(D) = w'(D) = w''(D) = 0$ and $w'(0) = 0$.
Hence four integrations by parts give
\[
W(z) = \frac{w_0}z + \frac{w''(0)}{z^3} + \frac{w'''(0) - e^{-zD}w'''(D)}{z^4} + \frac1{z^4}\int_0^De^{-zu}w''''(u)\,du,
\]
and the lemma follows.
The values of the derivatives are enclosed in ball arithmetic; the integral is decreasing in $\nu$, so it suffices to bound it at integers $\nu \in [30, 200]$.
\end{proof}

For example, $C_w(69.4, 69.4) = 51.80$, while \cite[Lemma~4]{BTY} uses the bound $52$ at this point.

\begin{lemma}\label{lem:L5}
Let $\sigma \in [0.993, 1)$ and $\eta \leq \eta_0 \leq 0.0145$.
Then $\Rp F(z) + \Rp F(2\sigma - 1 - \bar z) \geq 0$ for $0 \leq \Rp z \leq 2\sigma - 1$.
\end{lemma}

\begin{proof}
As in \cite[Lemma~5]{BTY}, it suffices to consider $\Rp z = 0$.
For fixed $\sigma$, the shifts in $F$ are exactly $(2\sigma-1)m$, and the substitution $y \mapsto y/(2\sigma - 1)$ reduces the claim to
\[
B(y) := \sum_{m=1}^7b_m\,\frac{m - \sgn(b_m)\eps_0/m}{m^2 + y^2} > 0 \quad (y \in \mathbb R), \qquad \eps_0 = \frac{C_w\eta_0^2}{w_0(2\sigma-1)^2},
\]
where $(1+x)\sum_m\kappa_mx^m = \sum_mb_mx^m$.
For $\eta_0 \leq 0.0145$ and $\sigma \geq 0.993$ we have $\eps_0 \leq 1/500$, using $C_w \leq 51.9$.
For $\eps_0 = 1/500$, the function $B$ is rational, its value at $0$ is positive, and its numerator has no real zeros, as we verify by exact root isolation.
Since $B$ decreases as $\eps_0$ increases, the claim follows.
\end{proof}

\begin{lemma}\label{lem:L89}
Let $x_0 = \sigma_0 - \frac12$, $C_1 = C_w(x_0/\eta_0, x_0/\eta_0)$, $C_2 = C_w(\sigma_0/\eta_0, \sigma_0/\eta_0)$ and $\Sigma_\kappa = \sum_m|\kappa_m|$.
Then the error terms $E_k$ of \cite[Lemma~8]{BTY} satisfy $\sum_ka_k|E_k| \leq E_2\lambda_0\eta^2 + E_3\eta^3$, where
\[
E_2 = a\Sigma_\kappa\frac{C_1}{\pi x_0^2}, \qquad E_3 = a_0\Sigma_\kappa\Big(\frac{23.6497\,C_1}{2\pi} + \frac{C_2}{x_0^3}\Big) + a\Sigma_\kappa\frac{C_2}{x_0^3}.
\]
\end{lemma}

\begin{proof}
We follow the proof of \cite[Lemma~8]{BTY} with two changes.
First, the inequality $\eta\log t \leq A_0$ used in \cite[(33)]{BTY} is replaced by the identity $\eta\log t = \lambda_0$.
Second, the integral \cite[(32)]{BTY} is evaluated with the digamma function itself: we have $\int_{\mathbb R}|\Rp\frac{\Gamma'}{\Gamma}(\frac14 + \frac{iy}2)|(x_0^2 + y^2)^{-3/2}\,dy \leq 23.6497$ for $x_0 \geq 0.493$, since the integral is decreasing in $x_0$.
\end{proof}

\begin{lemma}\label{lem:L1011}
\begin{enumerate}
\item[(i)] $f(u) \geq 0$ for $u \geq 0$.
\item[(ii)] $f(u) \geq \kappa f(0)$ for $0 \leq u \leq u_0(\eta_0, \sigma_0)$; for example, $u_0 = 59.15$, $49.46$ and $33.61$ for $\eta_0 = 0.0071$, $0.0085$ and $0.0125$, respectively.
\item[(iii)] $\sum_n\Lambda(n)n^{-\sigma}f(\log n)P(t\log n) \geq c_{10}\eta w_0$, where $c_{10} = \kappa^2\sum_{15\log p\leq u_0}m_p\log p$, and $m_p$ are the lower bounds for $\min_xG_p(1,x)$ in \cite[Table~1]{BTY}, except that $m_5 = 0.0073564$.
\end{enumerate}
\end{lemma}

\begin{proof}
(i) The function $w$ is non-negative on $[0, D]$; we verify this by bisection in ball arithmetic on $[0, D - 0.02]$ and by a Taylor bound near $D$, where $w'''(D) < 0$.
The polynomial $\sum_m\kappa_mx^m$ has no zero in $[0,1]$.

(ii) We verify $w(\eta_0u)\sum_m\kappa_me^{-d_mu} - \kappa w_0 \geq 0$, with $d_m$ as in \cite[Lemma~11]{BTY}, by bisection in ball arithmetic; near $u = 0$ we use that this function vanishes at $0$ and has derivative at least $1.19$ there.

(iii) As in \cite[Lemma~10]{BTY}, the terms with $n > e^{u_0}$, the primes with $15\log p > u_0$ and the prime powers $p^m$ with $m > 15$ may be omitted by (i) and $P \geq 0$.
The inequality $G_p(\sigma, x) \geq G_p(1, x)$ holds because $G_p(\sigma, x) = \sum_{m\leq15}p^{-m\sigma}P(mx)$ is non-increasing in $\sigma$.
The values $m_p$ are verified by bisection with a mean value enclosure; see Remark~\ref{rem:m5} for $p = 5$.
\end{proof}

\begin{lemma}\label{lem:L12}
In the notation of \cite[Lemma~12]{BTY}, we have
\[
\sum_ka_k\sum_m\kappa_m\Big(-\frac12\log\pi + \frac12\Rp\frac{\Gamma'}{\Gamma}\Big(\frac{s_k + \delta_m}{2} + 1\Big)\Big) \leq \frac{a\kappa}2\log t - G(\sigma_0),
\]
where $G(\sigma_0)$ is given by the formula in \cite{BTY}; in particular, $G(0.993) = 1.56839$ and $G(0.9935164) = 1.56888$.
\end{lemma}

Finally, the term $10^{-10}\eta$ of \cite[Lemma~13]{BTY} and the tail estimates in \cite[Lemma~6]{BTY} remain valid for $\eta \leq 0.015$ and $\log t \leq 45$; each tail term is at most $3\cdot10^{-9}$, so their sum is less than $2\cdot10^{-7}$.

\subsection{Evaluation of \texorpdfstring{$W$}{W}}

On $[0, D]$, the function $w$ is a linear combination of $u\cos(ku)$, $\cos(ku)$, $\sin(\phi - ku)$, $u$ and $1$, where $k = \tan\theta$.
Hence $W$ has a closed form.
With $J_0(s) = (1 - e^{-sD})/s$, $J_1(s) = (1 - e^{-sD}(1 + sD))/s^2$ and $\Sigma_\phi(z) = (e^{i\phi}J_0(z + ik) - e^{-i\phi}J_0(z - ik))/(2i)$, it reads
\begin{multline*}
W(z) = \sec^2\theta\Big[\frac{\sec^2\theta}{2}\sum_\pm\Big(\theta\cot\theta\,J_0 - \frac12J_1\Big)(z \mp ik)\\
+ (2\theta\cot\theta - 2)J_0(z) - J_1(z) + \frac{\Sigma_{2\theta}(z)}{\sin2\theta} - \frac{2\Sigma_\theta(z)}{\sin\theta}\Big].
\end{multline*}
In the verification we evaluate $W$ on complex balls with the second-order enclosure
\[
W(z) \in W(z_c) + W'(z_c)(z - z_c) + B\Big(0, \frac{r^2}2\sup|W''|\Big), \qquad r = |z - z_c|,
\]
where $W'$ has the same closed form with $J_0, J_1$ replaced by $-J_1, -J_2$.
Three integrations by parts of $u^2w(u)$ give $\sup|W''| \leq \min\{\int_0^Du^2we^{-\nu u},\ (2w_0 + \int_0^D|(u^2w)'''|e^{-\nu u})|z|^{-3}\}$.

\subsection{The zero-free region}

\begin{theorem}\label{thm:H}
Let $\lambda^*$ be the step function defined on the grid $l \in \{28, 28.25, \ldots, 43\}$ by the values computed in Section~\ref{sec:proof}, part of which are listed below, and extended by $\lambda^*(l) = \lambda^*(l')$ for $l \leq 43$, where $l'$ is the smallest grid point with $l' \geq l$, and by $\lambda^*(l) = 1/4.896$ for $l > 43$.
Then every zero $\beta + i\gamma$ of $\zeta(s)$ with $\gamma > H_0$ satisfies
\[
(1 - \beta)\log\gamma \geq \lambda^*(\log\gamma).
\]
\begin{center}\small
\begin{tabular}{c|cccccc}
\toprule
$l$ & 28 & 29 & 30 & 31 & 32 & 33\\
$\lambda^*(l)$ & 0.21424 & 0.21384 & 0.21345 & 0.21308 & 0.21272 & 0.21239\\
\midrule
$l$ & 34 & 35 & 36 & 37 & 38 & 40\\
$\lambda^*(l)$ & 0.21206 & 0.21175 & 0.21145 & 0.21117 & 0.21084 & 0.21033\\
\bottomrule
\end{tabular}
\end{center}
Moreover, $\lambda^*(l)/l$ is non-increasing.
\end{theorem}

\begin{proof}
For $\log\gamma > 43$ the claim is \cite[Theorem~1]{BTY}.
For $\log\gamma \leq 43$ we use an iteration with four steps, starting from $\lambda_0 = 1/4.896$.
At step $j$ we assume the conclusion of step $j-1$.
Since $\lambda_{j-1}(l)/l$ is non-increasing, this assumption provides the zero-free input in \eqref{eq:bty} with $A = \lambda_{j-1}(\log(Kt + T_0))$.
In every interval of heights $\log t \in (l - 0.25, l]$ we show $\mathcal M_1 > 0$ in ball arithmetic, for all $\log t$ in the interval and all $\lambda_0 = \eta\log t \in [\lambda_{j-1}(l), \lambda_j(l)]$, by adaptive bisection.
On each box, $\sigma$ is taken as a ball with $A$ at the top of the box, which is admissible since a smaller value of $A$ is a weaker hypothesis, and the constants are taken at $\eta_0 = \lceil\lambda_j(l)/l_{\min}\rceil_{10^{-4}} + 10^{-4}$ and at $\sigma_0$ rounded down.
Then a zero with $\lambda_{j-1}(l) \leq (1 - \beta)\log\gamma < \lambda_j(l)$ would satisfy \eqref{eq:bty} with $\mathcal M_1 > 0$, which is impossible.
The table of the fourth step differs from that of the third by at most $10^{-5}$ in each entry.
\end{proof}

For example, at $\log t = 29$ Theorem~\ref{thm:H} gives the zero-free region $\sigma > 1 - 1/(4.676\log t)$.

\section{Pairs of zeros close to the line \texorpdfstring{$\sigma = 1$}{sigma = 1}}\label{sec:rep}

In \cite[Lemma~6]{BTY}, the sum over zeros $\Rp\sum_\rho F(s_1 - \rho)$ is bounded from below by the contribution of the pair $\rho_0$, $1 - \bar\rho_0$, while every other zero $\rho$ is paired with $1 - \bar\rho$, and every such pair contributes a non-negative amount by Lemma~\ref{lem:L5}.
If there is a second zero close to $\rho_0$, we may keep its pair as well.

\begin{lemma}\label{lem:twozero}
Assume the situation of \eqref{eq:bty}, and let $\rho_1 = \beta_1 + i(t + d)$ be a second zero with $\beta_1 > \frac12$ and $\sigma - \beta_1 > 0$.
Let $\eta_1 = 1 - \beta_1$.
Then \eqref{eq:bty} holds with
\[
Z \geq a_1\Rp\big(F(\sigma - 1 + \eta_1 - id) + F(\sigma - \eta_1 - id)\big).
\]
\end{lemma}

\begin{proof}
The pair $\rho_1$, $1 - \bar\rho_1$ contributes exactly the displayed quantity, since the point $s_1$ enters with the coefficient $a_1$.
All other pairs contribute non-negative amounts by Lemma~\ref{lem:L5}, since $0 \leq \Rp(s_1 - \rho) \leq 2\sigma - 1$ for every zero below height $Kt + T_0$.
Since $F(\bar z) = \overline{F(z)}$, only $|d|$ matters.
\end{proof}

Consequently, two zeros at heights $t$ and $t + d$ cannot exist if
\[
\mathcal M_2 := \mathcal M_1 + a_1\Rp\big(F(\sigma - 1 + \eta_1 - id) + F(\sigma - \eta_1 - id)\big) > 0
\]
for one of the two choices of the primary zero, where $F$ is formed with the $\eta$ of the primary zero.
Heuristically, $\Rp F(\sigma - \eta_1 - id) \approx \eta w_0(\sigma - \eta_1)/((\sigma - \eta_1)^2 + d^2)$, which is about $0.04$ for $|d| \leq 0.05$ and $\eta, \eta_1 \approx 0.0075$; this term causes the repulsion.

\begin{theorem}\label{thm:rep}
Let $29 \leq l_k \leq 37$ be an integer, and let $k$ be an integer with $\lambda^*(l_k) < k\cdot10^{-4}l_k \leq 0.0125(l_k - 1)$.
There are explicit numbers $d^*(l_k, k) \geq 0$ with the following property.
If two zeros $\beta_j + i\gamma_j$ $(j = 0, 1)$ satisfy $H_0 < \gamma_j \leq e^{l_k}$ and $1 - \beta_j \leq k\cdot10^{-4}$, then $|\gamma_0 - \gamma_1| \geq d^*(l_k, k)$.
For each $l_k$, the value $d^*(l_k, k)$ is largest for the smallest admissible $k$ (for instance $0.142$ for $l_k = 29$), and for large $k$ it is at least $0.058$ for $l_k = 29$, $0.056$ for $l_k = 30$ and $0.044$ for $l_k = 37$.
Outside the stated range of $(l_k, k)$ we put $d^*(l_k, k) = 0$.
\end{theorem}

\begin{proof}
Let $e_{\min} = \lambda^*(l_k)/l_k$ and $e_{\max} = k\cdot10^{-4}$.
By Theorem~\ref{thm:H} and the monotonicity of $\lambda^*(l)/l$, both zeros satisfy $1 - \beta_j \geq e_{\min}$; since $e_{\min} > A/\log(Kt + T_0)$ (for example $0.00737 > 0.0067$ for $l_k = 29$), the hypothesis $\sigma - \beta_1 > 0$ of Lemma~\ref{lem:twozero} is satisfied.
Fix an interval of heights $(l - 1, l]$ with $l \leq l_k$.
We show $\mathcal M_2 > 0$ for all $(\log t, \eta_0, \eta_1, |d|) \in [l - 1, l + 10^{-12}] \times [e_{\min}, e_{\max}]^2 \times [0, d^*]$ in ball arithmetic by adaptive bisection, trying both orderings of the pair on each box; the extra width $10^{-12}$ accounts for the height of the second zero.
As in the proof of Theorem~\ref{thm:H}, $\sigma$ is a ball with $A$ at the top of the box, and the constants are taken at outward rounded values of $(\eta_0, \sigma_0)$.
We do this for the values of $k$ on a coarse grid, with $d^*$ chosen slightly below the largest value suggested by a floating-point computation; for the remaining $k$ we use the value of the next larger $k$ on the grid, which is admissible since enlarging $e_{\max}$ only enlarges the box.
Since a zero with height at most $e^{l_k}$ may lie in any lower interval of heights, we finally replace $d^*(l_k, k)$ by $\min_{l\leq l_k}d^*(l, k)$, where the minimum is taken over those $l$ for which $1 - \beta \leq k\cdot10^{-4}$ is compatible with Theorem~\ref{thm:H}.
\end{proof}

\section{Zero-density estimates}\label{sec:density}

In the computation, $N(s,t)$ is replaced by an upper bound $N^+(s,t)$, which is the minimum of the following bounds, each in its range of validity.
\begin{enumerate}
\item[(N1)] $N(0, T) \leq \frac{T}{2\pi}\log\frac{T}{2\pi e} + \log T + 1$ for $T \geq H_0$, which follows from \cite{HSW}.
\item[(N2)] The bounds of Kadiri, Lumley and Ng \cite[Table~1]{KLN}, in the form $N(\sigma, T) \leq A(\log T)^{5 - 2\sigma_0}T^{8(1-\sigma_0)/3} + B\log^2T$ for $\sigma \geq \sigma_0 \in \{0.60, 0.65, \ldots, 0.99\}$ and $T \geq 3.061\cdot10^{10}$.
\item[(N3)] The bounds $N(\sigma, T) \leq CT^{B(1-\sigma)}$ of Bellotti \cite[Theorem~1.1 and Table~2]{Bellotti} for $\sigma \geq 0.9927$ and $T > H_0$, both the uniform one and those for $e^{l-1} < T \leq e^l$ with $29 \leq l \leq 41$.
\item[(N4)] Theorem~\ref{thm:density} below, for $1 - \sigma \leq 0.03$ and $28 < \log T \leq 40$.
\end{enumerate}
In the binding range only (N4) matters.
In this section we prove Theorem~\ref{thm:density} by revisiting Bellotti's log-free zero-density estimate \cite[Sections~3--4]{Bellotti}.

\subsection{The log-free estimate}

In this section $L = \log T$, and we write $U = T^u$, $V = T^v$ with $v = \tau u$, $W = T^w$ and $X = T^x$.
Let $\psi_d$ be the Barban--Vehov weights, supported on $d \leq V$, and $\Psi(n) = \sum_{d\mid n}\psi_d$, so that $\Psi(1) = 1$ and $\Psi(n) = 0$ for $1 < n \leq U$.
Let $\theta_d$ be the Selberg-type weights, supported on $d \leq W$, and $\Theta(n) = \sum_{d\mid n}\theta_d$.
Let $\alpha = 1 - \delta$.
We divide the strip into rectangles of height $h = \kappa_0\delta$, choose in each rectangle one zero with $\beta \geq \alpha$ as its representative if there is one, and let $\tilde n$ be an upper bound for the number of zeros with $\beta \geq \alpha$ in a rectangle.

\begin{theorem}\label{thm:density}
Let $k \leq 300$ be a positive integer, $\delta_k = k\cdot10^{-4}$ and $\delta \in ((k-1)10^{-4}, \delta_k]$.
Let $l \in \{29, 29.5, 30, 30.5, 31, 32, \ldots, 40\}$ and $\log T \leq l$.
Then, for every choice of parameters $(u, \tau, w, x)$ which is admissible in the sense of Lemma~\ref{lem:D},
\[
N(1 - \delta, T) \leq N(1 - \delta_k, e^l) \leq \tilde n\,\frac{d_4d_5}{w}\cdot\frac{\mathcal C}{c_1^2 - c_5}\cdot\frac{e^{2x\delta_kl}}{\delta_kl},
\]
where the constants $\tilde n$, $d_4$, $d_5$, $\mathcal C$, $c_1$ and $c_5$ are given in Lemma~\ref{lem:D} with $\delta = \delta_k$ and $L = l$.
\end{theorem}

The first inequality in Theorem~\ref{thm:density} is the monotonicity of $N(\sigma, T)$; in particular, the bound is applied at the corner of each cell and no rescaling inside a cell is used.
For each cell $(k, l)$ we choose the parameters numerically so that the bound is small; since every admissible choice gives a valid bound, this optimisation needs no verification.

\begin{lemma}\label{lem:D}
Let $\alpha_k = 1 - \delta_k$ and $\nu = \lambda^*(l)/l$.
We call $(u, \tau, w, x)$ admissible if $\delta_k > \nu$ and
\begin{gather*}
0.02 < w < u, \qquad 1.3 \leq \tau \leq 2.5, \qquad uL \geq 40,\\
uL - 2\log L < xL, \qquad 1.1\,\tau u \leq x \leq 2.6\,\tau u, \qquad 2\nu xL \geq 1,
\end{gather*}
and if the constants $c_1$ and $c_5$ below satisfy $c_1 > 0$ and $c_1^2 - c_5 \geq 0.2$.
For admissible parameters, Theorem~\ref{thm:density} holds with the following constants.
\begin{enumerate}
\item[(i)] $c_1 = e^{-L^2/U} - B_1 - 3.8e^{-L(L-2u)}\dfrac{L^2 + 1}{w(v - u)L^6}$, where
\[
B_1 = \frac{10.45\cdot0.618\cdot5.922\,L}{2\pi}\,T^{(v+w)/2 + 1/6 + x(1/2 - \alpha_k)}.
\]
\item[(ii)] $d_5 = 1.0061\big(1 + (1.333 + 4e^{-wL/2})/(wL)\big)$ \cite[Lemma~4.3]{Bellotti}.
\item[(iii)] $d_4$ is any upper bound for $\sum_n\Psi(n)^2n^{-1}(e^{-n/X} - e^{-nL^2/U})$; we use Theorem~\ref{thm:G}.
\item[(iv)] $c_5 = \dfrac{d_4d_5}{2\pi}\cdot\dfrac{e^{wL}}{wL}\Big(e^{Lx(3/2 - 2\alpha_k)} + e^{(Lu - 2\log L)(3/2 - 2\alpha_k)}\Big)\cdot4.72\cdot0.618\,e^{L/6}L$.
\item[(v)] $\mathcal C = 1.0067\big(1 + 2\delta_k\Lambda(\kappa_0\delta_k)\big) + 2\sqrt{2\pi}\cdot2.4\cdot I'$, where $\kappa_0 \geq 1$ is arbitrary, $\Lambda$ is defined in Lemma~\ref{lem:pair} and $I'$ in Lemma~\ref{lem:crowd}.
\item[(vi)] $\tilde n = \max\{1, \lfloor M(\eta, l)/\mathrm{cost}(h/2)\rfloor\}$ for any $\eta > 0$, where $h = \kappa_0\delta_k$ and $M$ and $\mathrm{cost}$ are defined in \eqref{eq:Mbudget} and \eqref{eq:cost} below.
\end{enumerate}
\end{lemma}

\begin{proof}
We follow \cite{Bellotti} and only describe the changes.
The structure of the argument is that of \cite[(4.22)--(4.23)]{Bellotti}: if $J$ denotes the number of representatives, then
\[
J^2c_1^2 \leq J(c_2 + c_3 + c_4)\frac{X^{2-2\alpha}}{\delta\log T} + J^2c_5 \quad\text{and}\quad N(\alpha, T) \leq \tilde nJ.
\]
Each change below replaces one of Bellotti's constants by a smaller admissible one.

The term $B_1$ replaces $44.66\,d_2(v+w)^4T^{(v+w)/2 + x(1/2 - \alpha_0) + 27/164}\log^4T$ in \cite[(3.17)]{Bellotti}.
Writing $a = ga'$ and $b = gb'$ with $(a', b') = 1$ and using $|\psi_a|, |\theta_b| \leq 1$ and $\sum_{n\leq N}n^{-1/2} \leq 2\sqrt N$, we obtain the bound
\[
|M(\tfrac12 + it)| \leq \sum_gg^{-1/2}\cdot2\Big(\frac Vg\Big)^{1/2}\cdot2\Big(\frac Wg\Big)^{1/2} < 4\zeta(3/2)\sqrt{VW} < 10.45\sqrt{VW}
\]
for the mollifier, which contains no logarithms.
Moreover, we use $|\zeta(\frac12 + iu)| \leq 0.618|u|^{1/6}\log|u|$ for $3 \leq |u| \leq e^{105}$ by Hiary, Patel and Yang \cite{HPY}, $|\zeta(\frac12 + iu)| \leq 1.461$ for $|u| < 3$ \cite{Hiary}, and $(\gamma + |t|)^{1/6}\log(\gamma + |t|) \leq T^{1/6}\log T\,(1 + |t|)$.
This leads to the integral
\[
I_\Gamma = \sup_{-1/2\leq\sigma\leq1/2-\alpha}\int_{\mathbb R}|\Gamma(\sigma + it)|(1 + |t|)\,dt \leq 5.922 \qquad (\delta \leq 0.03),
\]
which we verify in Section~\ref{sec:proof}.
The last term of $c_1$ is Bellotti's bound for the tail $\sum_{n>U}\Psi(n)\Theta(n)n^{-\rho}e^{-nL^2/U}$ \cite[(3.19)ff.]{Bellotti}.

The last factor in $c_5$ replaces $J_1 + J_2 + J_3 + J_4 \leq 0.241 + 5.921 + 286.077T^{27/164} + 0.241$ in \cite[(4.17)--(4.19)]{Bellotti}.
For $\delta \leq 0.03$ we have $c := \frac32 - \beta_r - \beta_s \in [-\frac12, -0.44]$.
For all $|u| \geq 3$ we have $|\zeta(\frac12 + iu)| \leq \max\{1.461, 0.618|u|^{1/6}\log|u|\}$, since for $|u| > e^{105}$ the bound $66.7|u|^{27/164}$ of Patel and Yang \cite{PY} is smaller than the right-hand side.
Since $|u| \leq T + |t|$, the integral \cite[(4.16)]{Bellotti} is at most
\[
0.618\,T^{1/6}\log T\int_{\mathbb R}|\Gamma(c + it)|\Big(1 + \frac{|t|}3\Big)^{1/6}\Big(1 + \frac{\log(1 + |t|/3)}{\log3}\Big)dt \leq 4.72\cdot0.618\,T^{1/6}\log T;
\]
the integral is verified in Section~\ref{sec:proof} as well.

The quantity $d_4$ is the Barban--Vehov mean square of \cite[Lemma~4.2]{Bellotti}.
Since Theorem~\ref{thm:G} bounds the whole sum, Bellotti's condition $u + v < x$ is not needed.

For the off-diagonal terms $c_2 + c_3$, Bellotti assumes in \cite[(4.11)]{Bellotti} that representatives of distinct rectangles are at least $\delta$ apart.
This need not be the case, since representatives of adjacent rectangles may be arbitrarily close.
Lemma~\ref{lem:pair} below avoids this assumption, and it uses Theorem~\ref{thm:rep}.
The far term $c_4$ is treated by Lemma~\ref{lem:crowd}, and $\tilde n$ replaces $b_1\delta\log T + b_2$ in \cite[Lemma~2.4]{Bellotti}.
\end{proof}

\subsection{The pair sum}

Let $\mathcal Z(\eta) = -\zeta'(1 + \eta)/\zeta(1 + \eta)$.
For $s_0 = 1 + \eta + i\gamma_r$ with $\gamma_r \leq T$, the Hadamard product and Stirling's formula give
\begin{multline}\label{eq:Mbudget}
\sum_\rho\Rp\frac{1}{s_0 - \rho} \leq M(\eta, L)\\
:= \min\Big\{\frac12\log\frac T{2\pi} + \mathcal Z(\eta),\ \min_{1/2\leq c\leq2}\Big(\frac12\log\frac T{2\pi} + \frac{c^2 + 1/2}{2c}\mathcal Z(\eta) + \frac{\log(T/\pi)}{8c}\Big)\Big\},
\end{multline}
where the second expression comes from the non-negativity of $(c + \cos\varphi)^2$.
Every term in the sum is positive, and a zero with $\beta \geq 1 - \delta$ at vertical distance $d$ from $s_0$ contributes at least
\begin{equation}\label{eq:cost}
\mathrm{cost}(d) := \min\Big\{\frac{\eta}{\eta^2 + d^2},\ \frac{\eta + \delta}{(\eta + \delta)^2 + d^2}\Big\},
\end{equation}
since $a \mapsto a/(a^2 + d^2)$ is unimodal.

\begin{lemma}\label{lem:pair}
Fix a representative $r$ and let $h = \kappa_0\delta$.
On each side of $r$, the representative of the $j$-th rectangle has distance $d \in ((j-1)h, (j+1)h)$ from $r$.
Let $d^* = d^*(\lceil l\rceil, k)$ be as in Theorem~\ref{thm:rep} ($d^* = 0$ outside its range), and let
\[
b_j = \frac{g(m_j)}{m_j}, \qquad m_j = \max\{2\delta, (j-1)h, d^*\}, \qquad g(d) = \min\Big\{1, \frac{|\Gamma(1 + 2\delta + id)|}{\Gamma(1 + 2\delta)}\Big\}.
\]
Rectangles with $(j+1)h \leq d^*$ contain no representative.
Then
\[
\sum_{\substack{s\neq r\\ |\gamma_s - \gamma_r|\leq1}}|\Gamma(2 - \rho_r - \bar\rho_s)|X^{2-\beta_r-\beta_s} \leq \frac{1.0067\,X^{2\delta}}{2\delta}\cdot2\delta\Lambda(h),
\]
where $\Lambda(h)$ is the maximum of $\sum_jb_jy_j$ over $0 \leq y_j \leq 1$ (two variables for each $j$) subject to $\sum_j\mathrm{cost}((j+1)h)y_j \leq M(\eta, L) - \mathrm{cost}(0)$, for any $\eta > 0$.
\end{lemma}

\begin{proof}
We have $|\Gamma(z)| = |\Gamma(1 + z)|/|z|$, and for $0 \leq \Rp z = s \leq 2\delta$,
\[
|\Gamma(1 + s + id)| \leq \Gamma(1 + s)\frac{|\Gamma(1 + 2\delta + id)|}{\Gamma(1 + 2\delta)},
\]
since $\prod_n(1 + d^2/(1 + s + n)^2)^{-1}$ is increasing in $s$, and $\Gamma(1 + s) \leq 1$ for $0 \leq s \leq 1$.
Moreover, $X^y/y$ is increasing for $y \geq 1/\log X$, hence on $[2\nu, 2\delta]$ by the condition $2\nu xL \geq 1$; this gives the factor $1/m_j$, relative to the diagonal term $1.0067X^{2\delta}/(2\delta)$.
Applying \eqref{eq:Mbudget} at $s_0 = 1 + \eta + i\gamma_r$ bounds the total cost of the representatives near $\gamma_r$, and the zero $r$ itself costs at least $\mathrm{cost}(0)$.
Theorem~\ref{thm:rep} applies, since all representatives counted in the cell $(k, l)$ satisfy $1 - \beta \leq \delta_k$ and have height at most $e^l$; for half-integral $l$ we use the integer $\lceil l\rceil$.
\end{proof}

In the computation, $\Lambda(h)$ is bounded from above by the dual value $\mu(M(\eta, L) - \mathrm{cost}(0)) + \sum_j\max\{0, b_j - \mu\,\mathrm{cost}((j+1)h)\}$, which is valid for every $\mu \geq 0$.

\begin{lemma}\label{lem:crowd}
Let $\omega > 0$, $\eta > 0$ and $t \leq T$.
Then a window of height $\omega$ centred at $t$ contains at most
\[
n(\omega) = \min\Big\{\Big\lfloor\frac{\frac12\log\frac T{2\pi} + \mathcal Z(\eta)}{\min\{f_\omega(\eta), f_\omega(\eta + \delta)\}}\Big\rfloor,\ \Big\lceil\frac\omega\delta\Big\rceil + 1\Big\}, \qquad f_\omega(a) = \frac{a}{a^2 + \omega^2/4},
\]
representatives.
Consequently, Bellotti's constant $c_4$ \cite[(4.13)--(4.14)]{Bellotti} may be computed with
\[
I' = \delta\,n(\omega)\Big(\sum_{i\geq0}g_\alpha(1 + i\omega)\Big), \qquad g_\alpha(y) = e^{y/6}y^{3/2 - 2\alpha}e^{-\pi y/2}, \qquad \omega = \frac14,
\]
in place of $I = \int_1^Tg_\alpha(y)\,dy$.
\end{lemma}

\begin{proof}
The first claim follows from \eqref{eq:Mbudget} applied at $s = 1 + \eta + it$, as above.
For $y \geq 1$ the logarithmic derivative of $g_\alpha$ is $\frac16 + (\frac32 - 2\alpha)/y - \frac\pi2 < 0$, so $g_\alpha$ is decreasing and its maximum on the window $[1 + i\omega, 1 + (i+1)\omega]$ is attained at the left endpoint.
\end{proof}

\subsection{The Barban--Vehov mean square}

Let $\Delta = \log(V/U)$, $Y_1 = \log(X/V)$ and $Y_2 = \log(X/U)$, and let $\gamma_E$ denote Euler's constant.

\begin{theorem}\label{thm:G}
Let $\log U \geq 40$, $X > V$, $\tau \in [1.3, 2.5]$ and $\log X/\log V \in [1.1, 2.6]$.
Then
\[
\sum_n\frac{\Psi(n)^2}ne^{-n/X} \leq e^{-1/X} + D_4 + \mathrm{REM} + 2\sqrt{S_sS_t} + S_t,
\]
where
\[
D_4 = \frac{1}{\Delta^2}\Big[\frac{Y_2^2 - Y_1^2}2 - \gamma_E\Delta + k_1(Y_1 + Y_2 - 2\gamma_E) + 2k_0\Big], \qquad k_1 \leq -0.5853, \quad k_0 \leq 0.4214,
\]
$\mathrm{REM}$ is the explicit sum of contour remainders given in Lemma~\ref{lem:A-rem}, $S_t$ is the explicit bound of Lemma~\ref{lem:A-St} for the mean square of the non-smooth part of the weights, and $S_s = e^{-1/X} + D_4 + \mathrm{REM}$.
\end{theorem}

The proof is given in Appendix~\ref{app:G}.
In Lemma~\ref{lem:D} we take $d_4$ to be the right-hand side of Theorem~\ref{thm:G} minus $e^{-L^2/U}$; this is admissible since $\sum_n\Psi(n)^2n^{-1}e^{-nL^2/U} \geq e^{-L^2/U}$, the term $n = 1$ being $e^{-L^2/U}$.
The idea is to replace the weights by a Gaussian smoothing in $\log d$, whose Mellin transform decays like $e^{-\eta^2t^2/2}$; this allows us to close all contours inside a fixed box of height $400$, in which $1/\zeta$ is bounded by a rigorous computation, and the main term $e^{-1/X} + D_4$ arises from the residues at $s = 0$, $s = -w$ and $w = 0$.


For example, for the cell $(k, l) = (80, 30)$ of Theorem~\ref{thm:density}, at its parameters, $\mathrm{REM} \leq 0.016$, $S_t \leq 6.7\cdot10^{-5}$ and Theorem~\ref{thm:G} gives $d_4 \leq 1.62$.
Since $\log N^+$ depends on $d_4$ with elasticity about $1$, this is the most important improvement in the density estimate.
In the final tables, all cells have $\tau \in [1.54, 1.78]$, $\log X/\log V \in [1.42, 1.55]$ and $\log U \in [56.6, 78.6]$, and all of them use Theorem~\ref{thm:G}.

\section{Proof of Theorem~\ref{thm:main}}\label{sec:proof}

By Proposition~\ref{prop:transfer}, it suffices to cover $[43, 3159)$ by intervals $[L_1, L_2]$ and to establish \eqref{eq:psicert} on each of them with $q_{[L_1,L_2]} < 1$ whenever $L_2 > 57$.
We obtain \eqref{eq:psicert} from Proposition~\ref{prop:master}, in which $s_{\max} = 1 - \lambda^*(\log T_1)/\log T_1$ by Theorem~\ref{thm:H}, $S(s)$ is bounded by Lemma~\ref{lem:stieltjes}, and $N^+(s, t)$ is the minimum of the bounds (N1)--(N4) of Section~\ref{sec:density}.
For every interval $[L_1, L_2]$ we proceed as follows.
First we choose $c$ and a number $f \in [0.3, 0.98]$ by a search over a grid in floating-point arithmetic, and we choose $\eps$ by three steps of a fixed-point iteration so that the first term in \eqref{eq:master} is approximately $f\,\eps_R(L_2 + 1)/(1 + 2/L_1)$.
Then, with these values of $c$ and $\eps$ fixed, we evaluate the right-hand side of \eqref{eq:master} in ball arithmetic, using the actual value of the first term.
The integral over $s$ is bounded by $\sum_jS(s_j)\int_{s_j}^{s_{j+1}}e^{(s-1)L_1}\,ds$, which is admissible since $S$ is non-increasing; the grid consists of $61$ points up to $s_{\max} - 0.02$ followed by steps of length $1/(40L_1)$, and its last point is $s_{\max}$ rounded upwards.
In Lemma~\ref{lem:stieltjes} we use $80$ subintervals of equal length in $\log t$, the last point being $\log T_1$ rounded upwards, and a block is omitted only if the condition of Lemma~\ref{lem:stieltjes} is verified in ball arithmetic.
Finally, $q_{[L_1,L_2]}$ is defined as the upper endpoint of the resulting enclosure of $(1 + 2/L_1)\,|\psi(x) - x|/(x\,\eps_R^-(L_2+1))$, excluding $a(x)/x$.

We start with intervals of width $10$, and $2$ for $L_1 < 57$, and bisect an interval, down to width $0.5$, as long as its value of $q$ (or the floating-point value used in the search) is at least $1$.
For $L_1 \geq 1100$ the search for $(c, f)$ uses a coarser partition in Lemma~\ref{lem:stieltjes}; this affects only the choice of parameters.
For $L_1 < 57$ only some bound is needed, since this range enters only through the integral in \eqref{eq:pili}.
The results are as follows.
\begin{center}
\begin{tabular}{lll}
\toprule
range of $L$ & largest value of $q$ & remark\\
\midrule
$[43, 57)$ & $3.85$ on $[43, 43.5]$; at most $0.47$ on $[53, 57)$ & enters only through \eqref{eq:pili}\\
$[57, 1000)$ & $0.99137$ on $[872, 877]$ & width $5$ on $[857, 917)$\\
$[1000, 2000)$ & $0.8887$ on $[1000, 1010]$ & \\
$[2000, 2600)$ & $0.4064$ & \\
$[2600, 3159)$ & $0.4003$ & \\
\bottomrule
\end{tabular}
\end{center}
Hence $q_{[L_1, L_2]} < 1$ for every interval with $L_2 > 57$, and Proposition~\ref{prop:transfer} gives \eqref{eq:target}.
By Lemma~\ref{lem:reduce}, this proves the first statement of Theorem~\ref{thm:main}, and the second statement follows as explained in the introduction.
\qed

The binding range $L \approx 900$ corresponds to the bandwidth $T_1 \approx e^{30}$.
There the sum over zeros is dominated by hypothetical zeros with $(1 - \beta)\log\gamma \in [0.21, 0.32]$ and $\log\gamma \in (28.73, 30]$.
On the interval $[900, 910]$ the best kernel gives about $1.002$ (roughly $0.40$ for the smoothing term and $0.60$ for the zeros above $H_0$), whereas on intervals of width $5$ we obtain $q < 1$.

\subsection*{Verification}
All computations are performed in ball arithmetic with the Arb library \cite{arb} via \code{python-flint}.
Floating-point arithmetic enters only through the choice of parameters: the kernel parameters $(c, \eps)$, the parameters $(u, \tau, w, x)$ of Lemma~\ref{lem:D}, the auxiliary parameters $\eta$, $\kappa_0$ and $\mu$ in Lemmas~\ref{lem:pair} and~\ref{lem:crowd} and in Theorem~\ref{thm:G}, and the points of the partitions.
Each of them is fixed as an exact binary number before the bound is evaluated.
In detail, the computation consists of the following parts.
\begin{enumerate}
\item The enclosures of $\eps_R$ in Lemma~\ref{lem:epsR}, the constant in Lemma~\ref{lem:smooth}, and Proposition~\ref{prop:transfer}.
\item The constants of Section~\ref{sec:zfr}, the iteration in the proof of Theorem~\ref{thm:H}, and the cells of Theorem~\ref{thm:rep}.
\item The inputs of Theorem~\ref{thm:G}; in particular, the supremum of $|\mathcal A|$ over the torus of phases for $p \leq 100$ is evaluated on a grid with explicit cell enclosures, using interval arithmetic on double precision numbers with outward rounding after every operation.
\item The $3320$ non-empty cells of Theorem~\ref{thm:density}, evaluated from their stored parameters.
Here $\mathcal Z(\eta)$ is computed from the Taylor series of $\zeta$ at $1 + \eta$, $\Lambda(h)$ is bounded by its dual value, and $\tilde n$ and $n(\omega)$ are floors of verified upper bounds.
\item The two integrals in the proof of Lemma~\ref{lem:D}.
On $[-\frac12, -0.44] \times [0, 60]$ we verify $\Rp\frac{\Gamma'}{\Gamma}(\sigma + it) > 0$ on boxes, so that $|\Gamma(\sigma + it)|$ is increasing in $\sigma$ there; the integrals at the right endpoints are computed by rigorous numerical integration, and $|t| > 60$ is estimated explicitly.
This gives $I_\Gamma \leq 5.92164$ and $4.71363$ for the second integral.
\item The scan described above.
\end{enumerate}
On two processor cores the whole computation takes about $20$ hours.

\section{Remarks}\label{sec:remarks}

\subsection{On the literature}
While re-deriving the explicit results which we use, we noticed the following points.
None of them affects the main results of the papers concerned.

\begin{remark}\label{rem:btyf0}
In \cite{BTY}, the term coming from Lemma~10 enters the final inequality as $-0.1186\,w(0)\eta$, that is, with $f(0) = \eta w(0)$.
By \cite[Definition~2]{BTY}, however, $f(0) = \kappa\eta w(0)$ with $\kappa = 0.504$, so that this term should be $-0.0598\,w(0)\eta$.
Evaluating their inequality exactly at the critical height $t = e^{76.47}$ with this correction still gives $\eta\log t \geq 0.2054 > 1/4.896$, so \cite[Theorem~1]{BTY} is not affected.
We also note that $F(2\sigma - 1 + \beta_0)$ in \cite[(25)]{BTY} should read $F(\sigma - 1 + \beta_0)$, as in \cite[(47)]{BTY}.
\end{remark}

\begin{remark}\label{rem:m5}
For $p = 5$, \cite[Table~1]{BTY} gives $m_5 = 0.00735709$.
However, $G_5(1, x)$ takes the value $0.0073564919\ldots$ at $x = 1.6495242\ldots$.
All other entries of the table are confirmed.
The effect on the constant $0.23545$ in \cite{BTY} is $-10^{-6}$.
\end{remark}

\begin{remark}
The majorant $U(y)$ of \cite[(31)]{BTY} integrates to $28.3$ against $(x_0^2 + y^2)^{-3/2}$, and not to less than $24$ as stated in \cite[(32)]{BTY}.
With the digamma function itself the integral is $23.59$, so the inequality \cite[(32)]{BTY} is correct.
At $(\eta_0, \sigma_0) = (0.0071093, 0.9935164)$, Lemma~\ref{lem:L89} gives $185.4\eta^2 + 7972\eta^3$, compared with $187\eta^2 + 8013\eta^3$ in \cite[(48)]{BTY}.
\end{remark}

\begin{remark}
The uniform version of \cite[(21)]{BTY} replaces the shifts $(2\sigma - 1)m$ by $c_0m$ with $c_0 = 151/153$, which requires $\sigma \geq 0.993464$.
For $c_0 = 0.9864$, which corresponds to the smaller values of $\sigma$ used here, their test function takes negative values already for $\eps_0 = 1/2000$; Lemma~\ref{lem:L5} avoids this by working with a fixed $\sigma$.
Moreover, the numerical verification of $\partial_\sigma G_p < 0$ in \cite{BTY} is not needed, since it follows from $P \geq 0$, and the omission of the prime powers with $m > 15$ in \cite[Lemma~10]{BTY} is justified by $f \geq 0$.
\end{remark}

\begin{remark}\label{rem:buthe}
In \cite[Proposition~4]{Buthe2016}, quoted as \cite[Proposition~1]{Buthe2018}, the main term of the smoothing error is $2\eps x|\nu_c(\alpha)|\log(e^\eps x)/\log B$, where $\nu_c(t) = \int_{-\infty}^t\mu_c$ and $\mu_c(t) = -\int_{-\infty}^t\eta_{c,1}$ for $t < 0$.
With $\eta_{c,1}$ normalised as in Section~\ref{sec:explicit}, so that $\int\eta_{c,1} = 2$, one has $|\nu_c(0)| \sim \sqrt{2/(\pi c)}$, whereas \cite[Remark~1]{Buthe2018} gives $|\nu_c(0)| \sim (2\pi c)^{-1/2}$, which corresponds to a density of total mass $1$.
We have not been able to reconcile the two normalisations, and Lemma~\ref{lem:smooth} does not depend on this.
\end{remark}

\begin{remark}
In \cite[(4.11)]{Bellotti}, representatives of distinct rectangles are assumed to be at least $\delta$ apart, which does not seem to be justified; Lemma~\ref{lem:pair} avoids this assumption.
Since Bellotti's other constants are considerably larger than ours, her published bounds are not affected in practice.
\end{remark}

\subsection{On two recent preprints}

In \cite{De2}, it is claimed that \eqref{eq:ram} holds for every $x \geq e^{43.5102147}$.
The argument combines $\pi(x) < e\,\pi(x/e)$ with $\pi(x) > x/\log x$ in order to conclude in \cite[(29)]{De2} that $\pi(x)^2 - \frac{ex}{\log x}\pi(x/e) < 0$.
This conclusion requires $\frac{x}{\log x}e\pi(x/e) \geq \pi(x)^2$.
However, the first inequality only gives $\frac{x}{\log x}e\pi(x/e) > \frac{x}{\log x}\pi(x)$, and passing from this to $\pi(x)^2$ would require $\pi(x) \leq x/\log x$, which is the reverse of the second inequality and is false for $x \geq 17$.
The gap cannot be repaired by sharper inputs of the same kind, since the two sides of \eqref{eq:ram} differ by about $x^2/\log^6x$, while the bounds used there have relative accuracy of order $(\log x)^{-3/2}$.

In \cite{De1}, it is claimed that \eqref{eq:ram} holds for every $x \geq e^{547}$.
The monotonicity argument used there assumes Sterneck's conjecture $|M(x)| < \frac12\sqrt x$, which is false \cite{Neubauer}, and the paper asserts $M(x) = O(\sqrt x)$, which is an open problem; the proof given yields at most $O(x^{1/2+\eps})$.
Moreover, the main term retained there is itself only $O(x^2/\log^2x)$, so that the sign of $\pi(x)^2 - \frac{ex}{\log x}\pi(x/e)$ is not determined, and the range below $e^{3158.442}$ is treated by evaluating this function at $26$ sample points.
Hence neither preprint proves \eqref{eq:ram} in the range considered here.

\section*{Acknowledgement}
The computations and the drafting of this paper were carried out with substantial assistance from an AI system (Claude, Anthropic).
The arguments and the code were also reviewed by separate runs of the same system; these reviews are not a substitute for independent human verification.
All computations can be reproduced from the code in the repository mentioned in the introduction.

\appendix

\section{Proof of Theorem~\ref{thm:G}}\label{app:G}

Throughout this appendix the notation is local.
We fix the numerical parameters
\begin{gather*}
\eta = 0.02, \qquad T_2 = 400, \qquad W_0 = 40, \qquad c = 0.02, \qquad c' = 0.01,\\
\sigma_1 = 0.01, \qquad \theta_1 = 0.40, \qquad \theta_2 = 0.30, \qquad b = 0.38, \qquad \kappa_1 = 0.38,
\end{gather*}
where $\eta$ is the width of a Gaussian smoothing (not related to the $\eta$ of Sections~\ref{sec:zfr} and~\ref{sec:density}).
We write $U = e^{\ell_U}$, $V = e^{\ell_V}$ and $X = e^{\ell_X}$, so that $\Delta = \ell_V - \ell_U$, $Y_1 = \ell_X - \ell_V$ and $Y_2 = \ell_X - \ell_U$, and $E(s) = e^{\eta^2s^2/2}$.
We use the following explicit inputs.
\begin{enumerate}
\item[(B1)] In the box $\mathcal B = \{0.6 \leq \sigma \leq 1.3,\ |t| \leq 450\}$ we have $|1/\zeta(s)| \leq M_I := 14.57$, and $|\zeta(s)| \leq M_Z := 11.89$ if moreover $|s - 1| \geq 0.1$.
These bounds are obtained in ball arithmetic from Taylor models of the entire function $(s-1)\zeta(s)$ on a cover of $\mathcal B$ by small discs, with the remainder controlled by Cauchy's estimate and $|\zeta(s)| \leq |s|/|s-1| + |s|/\sigma$.
The function $1/\zeta$ is analytic in $\mathcal B$, since the Riemann hypothesis holds up to height $H_0$.
\item[(B2)] For $\sigma > 1$ and all $t$, $|\zeta(\sigma + it)| \leq \zeta(\sigma)$ and $|1/\zeta(\sigma + it)| \leq \zeta(\sigma)/\zeta(2\sigma)$.
For $0 < \sigma < 1$ and all $t$, $|\zeta(\sigma + it)| \leq (1/\sigma + 1/(1 - \sigma))(1 + |t|)$, which follows from $\zeta(s) = s/(s-1) - s\int_1^\infty\{u\}u^{-s-1}\,du$.
\item[(B3)] For $\sigma > -1$ the function $v \mapsto |\Gamma(\sigma + iv)|$ is decreasing on $[0, \infty)$, by the product formula for $|\Gamma|^2$.
\end{enumerate}

\subsection{Mellin representation}

Let $\rho(d) = 1$ for $d \leq U$, $\rho(d) = \log(V/d)/\Delta$ for $U < d \leq V$ and $\rho(d) = 0$ for $d > V$, so that $\psi_d = \mu(d)\rho(d)$ and $\Psi(n) = \sum_{d\mid n}\mu(d)\rho(d)$.
Since $\frac{1}{2\pi i}\int_{(\sigma)}z^sd^{-s}s^{-2}\,ds = \log^+(z/d)$ for $\sigma > 0$, we have
\[
\rho(d) = \frac1{2\pi i}\int_{(\sigma)}\hat\rho(s)d^{-s}\,ds, \qquad \hat\rho(s) = \frac{V^s - U^s}{\Delta s^2}.
\]
Let $Z$ be a standard normal random variable, and put $\rho_s(d) = \mathbb E[\rho(de^{-\eta Z})]$ and $\rho_t = \rho - \rho_s$.
Since $\mathbb E[e^{\eta Zs}] = E(s)$, the function $\rho_s$ has the Mellin transform $\hat\rho_s(s) = \hat\rho(s)E(s)$, which decays like a Gaussian on vertical lines: $|E(\sigma + it)| = e^{\eta^2(\sigma^2 - t^2)/2}$.
Let $\Psi_s(n) = \sum_{d\mid n}\mu(d)\rho_s(d)$, $\Psi_t = \Psi - \Psi_s$, and
\[
S_s = \sum_n\frac{\Psi_s(n)^2}ne^{-n/X}, \qquad S_t' = \sum_n\frac{\Psi_t(n)^2}ne^{-n/X}.
\]
By Minkowski's inequality in $\ell^2(e^{-n/X}/n)$,
\begin{equation}\label{eq:mink}
\sum_n\frac{\Psi(n)^2}ne^{-n/X} \leq \big(\sqrt{S_s} + \sqrt{S_t'}\big)^2 = S_s + 2\sqrt{S_sS_t'} + S_t'.
\end{equation}

\begin{lemma}\label{lem:A-euler}
For $\Rp w$, $\Rp(w + s_1)$, $\Rp(w + s_2)$ and $\Rp(w + s_1 + s_2)$ all positive,
\[
\sum_n\frac1{n^{1+w}}\sum_{d_1\mid n}\frac{\mu(d_1)}{d_1^{s_1}}\sum_{d_2\mid n}\frac{\mu(d_2)}{d_2^{s_2}} = \frac{\zeta(1 + w)\zeta(1 + w + s_1 + s_2)}{\zeta(1 + w + s_1)\zeta(1 + w + s_2)}\,\mathcal A(w, s_1, s_2),
\]
where $\mathcal A = \prod_p\mathcal A_p$ with
\[
\mathcal A_p = \frac{(1 - x_1 - x_2 + z)(1 - z)}{(1 - x_1)(1 - x_2)}, \qquad \mathcal A_p - 1 = \frac{zx_1 + zx_2 - z^2 - x_1x_2}{(1 - x_1)(1 - x_2)},
\]
$x_j = p^{-1-w-s_j}$ and $z = p^{-1-w-s_1-s_2}$.
The product converges absolutely if each sum of two of the real parts of $1 + w + s_1$, $1 + w + s_2$ and $1 + w + s_1 + s_2$, and twice the last of them, exceeds $1$, and $\mathcal A(w, 0, s_2) = \mathcal A(w, s_1, 0) = 1$.
\end{lemma}

\begin{proof}
With $q = p^{-1-w}$, the local factor of the left-hand side is $1 + \frac{q}{1-q}(1 - p^{-s_1})(1 - p^{-s_2}) = (1 - x_1 - x_2 + z)/(1 - q)$.
Dividing by the local factor $(1 - x_1)(1 - x_2)/((1 - q)(1 - z))$ of the zeta quotient gives $\mathcal A_p$.
For $s_1 = 0$ we have $x_1 = q$ and $z = x_2$, so $\mathcal A_p = 1$, and similarly for $s_2 = 0$.
The terms of $\mathcal A_p - 1$ are of order $p^{-2-\cdots}$, which gives absolute convergence.
\end{proof}

Since $\rho_s(d)$ decays faster than any power of $d$ as $d \to \infty$, Mellin inversion in $w$ (with $e^{-y} = \frac1{2\pi i}\int_{(c)}\Gamma(w)y^{-w}\,dw$) and in $s_1$, $s_2$ gives
\begin{equation}\label{eq:triple}
S_s = \frac1{(2\pi i)^3}\int_{(c)}\int_{(\sigma_1)}\int_{(\sigma_1)}\Gamma(w)X^w\hat\rho_s(s_1)\hat\rho_s(s_2)\,G(w, s_1, s_2)\,ds_2\,ds_1\,dw,
\end{equation}
where $G$ is the right-hand side of Lemma~\ref{lem:A-euler}; all integrals converge absolutely.

\subsection{The contour shifts}

\begin{proposition}\label{prop:A-main}
Let $K_c(w) = \zeta(1 + w)H_c(w)$, where
\[
H_c(w) = \frac1{2\pi i}\int_{\mathcal P}\frac{\mathcal A(w; -w-s, s)\,E(w + s)E(s)}{\zeta(1 - s)\zeta(1 + w + s)(w + s)^2s^2}\,ds,
\]
and $\mathcal P$ is the path $-c' - iT_2 \to \kappa_1 - iT_2 \to \kappa_1 + iT_2 \to -c' + iT_2$.
Then $H_c$ is analytic for $-\theta_2 \leq \Rp w \leq c$, $|\Ip w| \leq W_0$, and with $k_1 = H_c(0)$ and $k_0 = H_c'(0) + \gamma_EH_c(0)$ we have
\[
S_s \leq e^{-1/X} + D_4 + \mathrm{REM},
\]
where $D_4$ is as in Theorem~\ref{thm:G} and $\mathrm{REM}$ is the sum of the eight terms in Lemma~\ref{lem:A-rem}.
\end{proposition}

\begin{proof}
We write \eqref{eq:triple} as $S_s = \frac1{2\pi i}\int_{(c)}\Gamma(w)X^wF_s(w)\,dw$ with $F_s(w) = \sum_n\Psi_s(n)^2n^{-1-w}$, and split the $w$-integral into $|\Ip w| \leq W_0$ and $|\Ip w| > W_0$.
On the first part we move the contours in four steps.

\emph{Step 1.}
We move $s_2$ to $\Rp s_2 = -c'$.
The only pole crossed is $s_2 = 0$, where $\hat\rho_s$ has a simple pole with residue $1$; since $G(w, s_1, 0) = 1$ and $\frac1{2\pi i}\int\hat\rho_s(s_1)\,ds_1 = \rho_s(1)$, the residue contributes $\rho_s(1)\cdot\frac1{2\pi i}\int_{|\Ip w|\leq W_0}\Gamma(w)X^w\,dw$.
Adding the part $|\Ip w| > W_0$ of $S_s$, we obtain
\[
\rho_s(1)e^{-1/X} + \frac1{2\pi i}\int_{|\Ip w| > W_0}\Gamma(w)X^w\big(F_s(w) - \rho_s(1)\big)\,dw.
\]
Since $\Psi_s(1) = \rho_s(1) \in [0, 1]$ and $|\Psi_s(n)| \leq \tau(n)$, we have $|F_s(w) - \rho_s(1)| \leq \frac14 + \sum_{n\geq2}\tau(n)^2n^{-1-c} \leq \zeta(1+c)^4/\zeta(2 + 2c)$, so the second term is at most the first term of $R_8$ in Lemma~\ref{lem:A-rem}.
The pole of $\zeta(1 + w + s_1 + s_2)$ at $s_2 = -w - s_1$ has real part $-c - \sigma_1 < -c'$ and is not crossed.

\emph{Step 2.}
For $|\Ip s_1| \leq T_2$ we move $s_1$ to $\Rp s_1 = -\theta_1$, joining the moved segment to the original line $\Rp s_1 = \sigma_1$ by horizontal connectors at $\Ip s_1 = \pm T_2$.
In the region swept, $1 + w + s_1 \in \mathcal B$, so $1/\zeta(1 + w + s_1)$ is analytic.
At $s_1 = 0$ the function $\hat\rho_s$ has a simple pole with residue $1$ and $G(w, 0, s_2) = 1$, so this residue equals
\[
\frac{1}{2\pi i}\int_{(-c')}\hat\rho_s(s_2)\,ds_2\cdot\frac1{2\pi i}\int_{|\Ip w|\leq W_0}\Gamma(w)X^w\,dw = (\rho_s(1) - 1)\big(e^{-1/X} - \vartheta\big),
\]
where $|\vartheta|$ is at most the first term of $R_8$; since $\rho_s(1) \leq 1$ and $e^{-1/X} - \vartheta > 0$, this residue is non-positive and may be omitted.
The pole at $s_1 = -w - s_2$, which is crossed when $|\Ip(w + s_2)| \leq T_2$, gives
\[
R_x = \frac1{(2\pi i)^2}\iint\Gamma(w)X^w\hat\rho_s(-w-s_2)\hat\rho_s(s_2)\,\frac{\zeta(1+w)\,\mathcal A(w; -w-s_2, s_2)}{\zeta(1 - s_2)\zeta(1 + w + s_2)}\,ds_2\,dw.
\]
The integral over the moved segment is $R_1$, the integrals over the connectors and over the remaining parts of $\Rp s_1 = \sigma_1$ make up $R_2$, and extending $R_x$ to all $s_2$ costs $R_3$; the integrand of $R_2$ has an integrable singularity at $s_1 = -w - s_2$, and Fubini's theorem applies.

\emph{Step 3.}
We expand
\[
\Delta^2\hat\rho_s(-w-s_2)\hat\rho_s(s_2)(w + s_2)^2s_2^2 = \Big(V^{-w} - V^{-w}\Big(\frac UV\Big)^{s_2} - U^{-w}\Big(\frac VU\Big)^{s_2} + U^{-w}\Big)E(w + s_2)E(s_2),
\]
and denote by $\Phi(s_2)$ the remaining factor of the integrand.
Then $\Phi$ has a simple pole at $s_2 = 0$ with residue $-E(w)/w^2$ (from the zero of $1/\zeta(1 - s_2)$) and a simple pole at $s_2 = -w$ with residue $E(w)/w^2$ (from the zero of $1/\zeta(1 + w + s_2)$), since $\mathcal A(w; -w, 0) = \mathcal A(w; 0, -w) = 1$.
In each of the four terms we keep the parts of the line $\Rp s_2 = -c'$ with $|\Ip s_2| > T_2$; together they give $R_4$.
On $|\Ip s_2| \leq T_2$ we proceed as follows, always joining the moved segment to the remaining line by horizontal connectors at $\Ip s_2 = \pm T_2$.
For the two diagonal terms we move the segment to $\mathcal P$; this crosses $s_2 = 0$ and gives $(V^{-w} + U^{-w})(K_c(w) + E(w)/w^2)$.
In the term with $(U/V)^{s_2}$ we move the segment to $\Rp s_2 = b$, crossing $s_2 = 0$; this gives $-V^{-w}E(w)/w^2$, and the moved segment and its connectors give $R_5$.
In the term with $(V/U)^{s_2}$ we move the segment to $\Rp s_2 = -b$, crossing $s_2 = -w$; this gives $-U^{-w}(V/U)^{-w}E(w)/w^2 = -V^{-w}E(w)/w^2$, and the moved segment and its connectors give $R_6$.
In all these regions $1 - s_2$ and $1 + w + s_2$ have imaginary parts of absolute value at most $T_2 + W_0 = 440$, so (B1) applies where the real parts are at most $1.3$.
Altogether,
\begin{multline*}
R_x = \frac1{\Delta^2}\cdot\frac1{2\pi i}\int\Gamma(w)X^w\Big[V^{-w}\Big(K_c(w) - \frac{E(w)}{w^2}\Big) + U^{-w}\Big(K_c(w) + \frac{E(w)}{w^2}\Big)\Big]dw\\
+ O^*(R_3 + R_4 + R_5 + R_6),
\end{multline*}
where $O^*(R)$ denotes a quantity of absolute value at most $R$.

\emph{Step 4.}
Finally we move $w$ to $\Rp w = -\theta_2$ for $|\Ip w| \leq W_0$.
The integrand has a pole only at $w = 0$, since $X^wV^{-w} = e^{Y_1w}$ and $X^wU^{-w} = e^{Y_2w}$.
Write $\Gamma(w)e^{Yw} = w^{-1} + (Y - \gamma_E) + (\frac{Y^2}2 - \gamma_EY + \frac{\gamma_E^2}2 + \frac{\pi^2}{12})w + O(w^2)$, $K_c(w) = k_1w^{-1} + k_0 + O(w)$ and $E(w)/w^2 = w^{-2} + \eta^2/2 + O(w^2)$.
The residue of $\Gamma(w)e^{Yw}K_c(w)$ is $k_1(Y - \gamma_E) + k_0$, and that of $\Gamma(w)e^{Yw}E(w)/w^2$ is $\frac{Y^2}2 - \gamma_EY + \frac{\gamma_E^2}2 + \frac{\pi^2}{12} + \frac{\eta^2}2$.
Hence the residues at $w = 0$ sum to
\[
\frac{1}{\Delta^2}\Big[\frac{Y_2^2 - Y_1^2}2 - \gamma_E(Y_2 - Y_1) + k_1(Y_1 + Y_2 - 2\gamma_E) + 2k_0\Big] = D_4,
\]
since $Y_2 - Y_1 = \Delta$.
The integral over $\Rp w = -\theta_2$ is $R_7$, and the connectors at $|\Ip w| = W_0$ give the second term of $R_8$.
\end{proof}

\subsection{The remainders}

For a real number $\sigma$ and $0 \leq t_0 < t_1 \leq \infty$ put
\[
\mathcal J_\nu(\sigma; t_0, t_1) = 2\int_{t_0}^{t_1}e^{\eta^2(\sigma^2 - t^2)/2}\frac{(1 + t)^\nu}{\sigma^2 + t^2}\,dt, \qquad h(\sigma) = \frac{V^\sigma + U^\sigma}{\Delta},
\]
so that $\int_{t_0\leq|t|\leq t_1}|\hat\rho_s(\sigma + it)|(1 + |t|)^\nu\,dt \leq h(\sigma)\mathcal J_\nu(\sigma; t_0, t_1)$.
Let $\mathcal G_\nu = \int_{|v|\leq W_0}|\Gamma(c + iv)|(1 + |v|)^\nu\,dv$.
For sets $I_1, I_2, I_3$ of real numbers let $\mathcal A^*(I_1, I_2, I_3)$ denote an upper bound for $|\mathcal A(w, s_1, s_2)|$ whenever the real parts of $1 + w + s_1$, $1 + w + s_2$ and $1 + w + s_1 + s_2$ lie in $I_1$, $I_2$ and $I_3$ (a single number stands for a one-point set).
Let $Z_b = \max\{M_Z, 1/\sigma_2 + 1/(1 - \sigma_2)\}$ with $\sigma_2 = 1 + c - \theta_1 - c'$, $\epsilon_T = e^{-\eta^2T_2^2/2}$, $\epsilon_T' = e^{-\eta^2(T_2 - W_0)^2/2}$, $\zeta_1 = \zeta(1 + c)$ and $\iota(\sigma) = \zeta(\sigma)/\zeta(2\sigma)$.

\begin{lemma}\label{lem:A-rem}
Let $\mathcal A_1 = \mathcal A^*(1 + c - \theta_1, 1 + c - c', \sigma_2)$, $\mathcal A_1^\square = \mathcal A^*([1 + c - \theta_1, 1 + c + \sigma_1], 1 + c - c', [\sigma_2, 1 + c + \sigma_1 - c'])$, $\mathcal A_x = \mathcal A^*(1 + c', 1 + c - c', 1)$, $\mathcal A_V = \mathcal A^*(1 - b, 1 + c + b, 1)$ and $\mathcal A_U = \mathcal A^*(1 + b, 1 + c - b, 1)$.
Then the remainders in the proof of Proposition~\ref{prop:A-main} satisfy
\begin{align*}
R_1 &\leq \frac{\mathcal G_1e^{c\ell_X}\zeta_1Z_bM_I\,\iota(1 + c - c')\mathcal A_1}{(2\pi)^3}\,h(-\theta_1)\mathcal J_1(-\theta_1; 0, T_2)\,h(-c')\mathcal J_1(-c'; 0, \infty),\\
R_2 &\leq \frac{\zeta_1\iota(1 + c - c')\mathcal A_1^\square}{(2\pi)^3}\Big[2\mathcal G_0e^{c\ell_X}M_I\,\hbar\,\mathcal Z\\
&\qquad + \mathcal G_1e^{c\ell_X}h(\sigma_1)\mathcal J_1(\sigma_1; T_2, \infty)\,\zeta(1 + c + \sigma_1 - c')\,\iota(1 + c + \sigma_1)\,h(-c')\mathcal J_1(-c'; 0, \infty)\Big],\\
R_3 &\leq \frac{\mathcal G_0e^{c\ell_X}\zeta_1\mathcal A_x\iota(1 + c')\iota(1 + c - c')}{(2\pi)^2}\cdot\frac{(U^{-(c-c')} + V^{-(c-c')})e^{\eta^2(c-c')^2/2}\epsilon_T}{\Delta T_2^2}\,h(-c')\mathcal J_0(-c'; 0, \infty),\\
R_4 &\leq \frac{4\mathcal G_0e^{cY_2}\zeta_1\mathcal A_x\iota(1 + c')\iota(1 + c - c')}{(2\pi)^2\Delta^2}\cdot\frac{2e^{\eta^2((c-c')^2 + c'^2)/2}\epsilon_T}{3(T_2 - W_0)^3},\\
R_5 &\leq \frac{\mathcal G_0e^{cY_1}}{\Delta^2}\big(r_V + r_V'\big), \qquad
R_6 \leq \frac{\mathcal G_0e^{cY_2}}{\Delta^2}\big(r_U + r_U'\big), \qquad
R_7 \leq \frac{(e^{-\theta_2Y_1} + e^{-\theta_2Y_2})\,\mathcal R}{2\pi\Delta^2},\\
R_8 &\leq \frac{\gamma_0e^{c\ell_X}\zeta_1^4}{2\pi\zeta(2 + 2c)}\\
&\qquad + \frac{2\gamma_0(\theta_2 + c)e^{cY_2}\mathcal K}{\Delta^2},\\
\gamma_0 &= \frac{4\cdot3.4\,W_0^{1/2}e^{-\pi W_0/2}}{\pi(1 - 1/(\pi W_0))},
\end{align*}
where
\begin{align*}
r_V &= \frac{e^{-\Delta b}\zeta_1\mathcal A_VM_I\iota(1 + c + b)}{2\pi}\cdot\frac{\pi e^{\eta^2((c+b)^2 + b^2)/2}}{b^2(b + c)},\\
r_U &= \frac{e^{-\Delta b}\zeta_1\mathcal A_UM_I\iota(1 + b)}{2\pi}\cdot\frac{\pi e^{\eta^2((c-b)^2 + b^2)/2}}{b^2(b - c)},\\
r_V' &= \frac{2(b + c')e^{c'\Delta}e^{\eta^2((c+b)^2 + b^2)/2}\epsilon_T\epsilon_T'}{2\pi T_2^2(T_2 - W_0)^2}\\
&\qquad\times\zeta_1\,\mathcal A^*([1 - b, 1 + c'], [1 + c - c', 1 + c + b], 1)\,M_I\,\iota(1 + c - c'),\\
r_U' &= \frac{2(b - c')e^{\eta^2b^2}\epsilon_T\epsilon_T'}{2\pi T_2^2(T_2 - W_0)^2}\\
&\qquad\times\zeta_1\,\mathcal A^*([1 + c', 1 + b], [1 + c - b, 1 + c - c'], 1)\,\iota(1 + c')\,M_I.
\end{align*}
Here $\hbar = (V^{\sigma_1} + U^{\sigma_1})e^{\eta^2(\theta_1^2 - T_2^2)/2}/(\Delta T_2^2)$ bounds $|\hat\rho_s|$ on the connectors of Step~2,
\begin{multline*}
\mathcal Z = (\theta_1 + \sigma_1)\Big(\Big(2 + \frac{2 + W_0 + T_2}{\sigma_2}\Big)h(-c')\mathcal J_0(-c'; 0, \infty) + \frac{h(-c')\mathcal J_1(-c'; 0, \infty)}{\sigma_2}\Big)\\
+ 2\pi\sqrt{(\theta_1 + \sigma_1)^2 + 1}\,\frac{(V^{-c'} + U^{-c'})e^{\eta^2c'^2/2}}{\Delta c'^2},
\end{multline*}
$\mathcal R = 2\int_0^{W_0}|\Gamma(-\theta_2 + iv)|\big(\mathcal K(v) + e^{\eta^2\theta_2^2/2}\theta_2^{-2}\big)\,dv$, $\mathcal K = \mathcal K(W_0) + e^{\eta^2\theta_2^2/2}W_0^{-2}$, and
\begin{multline*}
\mathcal K(v) = \frac{M_Z}{2\pi}\Big(\mathcal A^*(1 - \kappa_1, [1 + \kappa_1 - \theta_2, 1 + \kappa_1 + c], 1)M_I\,\iota(1 + \kappa_1 - \theta_2)\frac{\pi(2\kappa_1 - \theta_2)e^{\eta^2((\kappa_1 + c)^2 + \kappa_1^2)/2}}{\kappa_1(\kappa_1 - \theta_2)((2\kappa_1 - \theta_2)^2 + v^2)}\\
+ \mathcal A^*([1 - \kappa_1, 1 + c'], [1 - \theta_2 - c', 1 + c + \kappa_1], 1)M_I^2\,\frac{2(\kappa_1 + c')e^{\eta^2((\kappa_1 + c)^2 + \kappa_1^2)/2}\epsilon_T\epsilon_T'}{T_2^2(T_2 - W_0)^2}\Big).
\end{multline*}
The function $\mathcal K(v)$ bounds $|K_c(w)|$ for $\Rp w = -\theta_2$ and $|\Ip w| = v$, and also for $-\theta_2 \leq \Rp w \leq c$ and $|\Ip w| = W_0$.
\end{lemma}

\begin{proof}
Each term is the integral of the absolute value of the integrand over the corresponding piece of contour, bounded factor by factor; several of the resulting bounds are weaker than necessary by small constant factors.
The factor $\Gamma(w)X^w$ gives $\mathcal G_\nu e^{c\ell_X}$, or $\mathcal G_\nu e^{cY_j}$ after $V^{-w}$ or $U^{-w}$ has been absorbed; in $R_4$, all four coefficients are at most $e^{cY_2}$, since $e^{cY_1 + c'\Delta} \leq e^{cY_2}$.
The factors $\hat\rho_s$ give $h(\sigma)\mathcal J_\nu(\sigma; \cdot, \cdot)$.
The zeta factors are bounded by (B1) when the argument lies in $\mathcal B$, and by (B2) otherwise; for real parts in $[1.3, 1.4]$ we use $\iota(1.3) < M_I$.
In $R_1$ we use $|\zeta(1 + w + s_1 + s_2)| \leq Z_b(1 + |\Ip w|)(1 + |\Ip s_1|)(1 + |\Ip s_2|)$ on $\Rp(1 + w + s_1 + s_2) = \sigma_2$, which follows from (B1) for heights at most $450$ and from (B2) beyond; the three factors $1 + |\Ip|$ are absorbed into $\mathcal G_1$ and $\mathcal J_1$.
In $R_2$, on the connectors, the pole of $\zeta(1 + u)$ with $u = w + s_1 + s_2$ is handled by $|\zeta(1 + u)| \leq 1/|u| + 1 + |1 + u|/\Rp(1 + u)$, integrating $1/|u|$ over a disc of radius $\sqrt{(\theta_1 + \sigma_1)^2 + 1}$ and using $1/|u| \leq 1$ elsewhere.
Whenever a contour runs at height at least $T_2$, the factor $|E(s)|$ gives $\epsilon_T$ and, where it is present, $|E(w + s)|$ gives $\epsilon_T'$, while $|s| \geq T_2$ and $|w + s| \geq T_2 - W_0$.
The identity $\int_{\mathbb R}dt/((t^2 + a^2)((t + v)^2 + a'^2)) = \pi(a + a')/(aa'((a + a')^2 + v^2))$ gives the factors in $r_U$, $r_V$ and $\mathcal K(v)$.
In $\mathcal K(v)$, $|\zeta(1 + w)| \leq M_Z$ holds since $|w| \geq \theta_2$ on $\Rp w = -\theta_2$ and $|w| \geq W_0$ on the connectors.
In $R_8$ we use the bound for $|F_s(w) - \rho_s(1)|$ from Step~1 and, for $v \geq W_0$ and $-\theta_2 \leq \sigma \leq c$,
\begin{multline*}
|\Gamma(\sigma + iv)| = \frac{|\Gamma(\sigma + 1 + iv)|}{|\sigma + iv|} \leq \Gamma(\sigma + 1)\frac{|\Gamma(2 + iv)|}{v}\\
= \Gamma(\sigma + 1)\frac{\sqrt{1 + v^2}}{v}\sqrt{\frac{\pi v}{\sinh\pi v}} \leq 3.4\,v^{1/2}e^{-\pi v/2},
\end{multline*}
where we used that $|\Gamma(x + iv)|/\Gamma(x)$ is increasing in $x > 0$ and $\Gamma(\sigma + 1) \leq \Gamma(0.7) < 1.3$.
\end{proof}

All quantities in Lemma~\ref{lem:A-rem} which do not depend on $(U, V, X)$ (the integrals $\mathcal G_\nu$, $\mathcal J_\nu$ and $\mathcal R$, the values of $\zeta$, and the bounds $\mathcal A^*$) are computed once in ball arithmetic.
The integrals are upper Riemann sums, using (B3) and the monotonicity of the other factors.
For $\mathcal A^*$ we bound $\prod_{p\leq100}|\mathcal A_p|$ by the product of the suprema of $|\mathcal A_p|$ over the $3$-torus of phases of $x_1$, $x_2$ and $z$, evaluated on a grid of $96^3$ cells with an explicit cell enclosure (moving a phase by at most $h/2$ moves $x_j$ by at most $|x_j|h/2$ and $z$ by at most $3|z|h/2$), and with the moduli of $x_1$, $x_2$, $z$ given as intervals when the real parts vary.
Where the real parts vary along a contour, the range is cut into up to $32$ pieces, according to the real parts of the moving variables, and the maximum over the pieces is taken.
For $100 < p \leq 10^6$ we use $\log|\mathcal A_p| \leq |\mathcal A_p - 1|$ and the triangle inequality, and for $p > 10^6$ the bound $4.2\cdot10^{6(1-m)}/(m - 1)$, where $m$ is the least of the four sums of pairs of real parts in Lemma~\ref{lem:A-euler}.
For example, $\mathcal A_1 \leq 1202.27$ and $\mathcal A_x \leq 14.56$.

\subsection{The non-smooth part}

\begin{lemma}\label{lem:A-bsum}
Let $\varpi(y) = \phi(|y|) - |y|\bar\Phi(|y|)$, where $\phi$ and $\bar\Phi$ are the standard normal density and tail, and $\Psi_1(y) = \int_y^\infty\varpi$.
For real $A$ put $b(A) = \sum_{r\geq1}\varpi((\log r - A)/\eta)/r$.
Then
\[
b(A) \leq \eta\Big(\frac12 - \Psi_1\Big(\frac A\eta\Big)\Big) + 0.8\,e^{\eta^2/2}e^{-A} \quad (A \geq 0), \qquad b(A) \leq \eta\Psi_1\Big(\frac{|A|}\eta\Big) + 2\varpi\Big(\frac{|A|}\eta\Big) \quad (A < 0).
\]
In particular $b(A) \leq \eta/2 + 0.8e^{\eta^2/2}$ for all $A$, and $b(A) \leq \eta/2 + 0.8e^{\eta^2/2}e^{-A}$ for $A \geq 0$.
\end{lemma}

\begin{proof}
The function $\varpi$ is even and log-concave, since $\varpi(y) = \int_{|y|}^\infty\bar\Phi$ for $y \geq 0$ and $\bar\Phi$ is log-concave.
Hence $u \mapsto \varpi((u - A)/\eta)e^{-u}$ is log-concave, and $g(t) = \varpi((\log t - A)/\eta)/t$ is unimodal on $[1, \infty)$.
For a unimodal function, $\sum_{r\geq1}g(r) \leq \int_1^\infty g + 2\max g$.
Here $\int_1^\infty g(t)\,dt = \eta\int_{-A/\eta}^\infty\varpi = \eta\Psi_1(-A/\eta)$, and $\Psi_1(-y) = \frac12 - \Psi_1(y)$ since $\int_{\mathbb R}\varpi = \frac12$.
Since $\varpi \leq \phi$ and $\phi(y)e^{-\eta y} = \phi(y + \eta)e^{\eta^2/2}$, we have $\max g \leq e^{-A}\phi(0)e^{\eta^2/2} \leq 0.4e^{\eta^2/2}e^{-A}$.
For $A < 0$ the function $g$ is decreasing on $[1, \infty)$, so $\sum_rg(r) \leq g(1) + \int_1^\infty g = \varpi(|A|/\eta) + \eta\Psi_1(|A|/\eta)$.
\end{proof}

\begin{lemma}\label{lem:A-St}
We have $|\rho_t(d)| \leq \frac\eta\Delta[\varpi(\frac{\log d - \ell_U}\eta) + \varpi(\frac{\log d - \ell_V}\eta)]$ for every $d \geq 1$.
Let $\beta(m) = b(\ell_U - \log m) + b(\ell_V - \log m)$, and let $\Sigma$ be an upper bound for $\sum_m\mu^2(m)\varphi(m)^{-1}\beta(m)^2$.
Then, for every $\epsilon > 0$,
\[
S_t' \leq S_t := \Big(\frac\epsilon e\Big)^\epsilon X^\epsilon\Big(\frac1\epsilon + 1\Big)\Big(\frac\eta\Delta\Big)^2\Sigma + \bar\Phi(\ell_U/\eta)^2,
\]
and an admissible value of $\Sigma$ is
\[
\Sigma = \beta_0^2\Big(\log10^6 + 1.3326 + \frac4{1000}\Big) + \sum_j\beta_j^2\Big(\delta_u + 8e^{-u_j/2}\Big) + \frac{36(\eta + 2)^2e^{-y_0^2}}{2\pi}\Big(1 + \frac{\eta}{2y_0}\Big),
\]
where $\beta_0 = \eta + 0.8e^{\eta^2/2}(e^{-(\ell_U - \log10^6)} + e^{-(\ell_V - \log10^6)})$, the $u_j = \log10^6 + j\delta_u$ with $\delta_u = 0.02$ run over $[\log10^6, \ell_V + 60\eta)$, $\beta_j$ is the upper bound for $\beta(m)$ on $\log m \in [u_j, u_j + \delta_u]$ given by Lemma~\ref{lem:A-bsum} (using the monotonicity of the bounds in $A$), and $y_0 = 60$.
\end{lemma}

\begin{proof}
As a function of $y = \log d$, $\rho$ is piecewise linear with a change of slope of $-1/\Delta$ at $\ell_U$ and of $+1/\Delta$ at $\ell_V$.
Since $\mathbb E[(y - \eta Z)^+] - y^+ = \eta\varpi(y/\eta)$, this gives the bound for $\rho_t$.
Next, $e^{-y} \leq (\epsilon/e)^\epsilon y^{-\epsilon}$ for $y > 0$, and for any $f$,
\[
\sum_n\frac1{n^{1+\epsilon}}\Big(\sum_{d\mid n}\mu(d)f(d)\Big)^2 = \zeta(1 + \epsilon)\sum_m\frac{\mu^2(m)}{m^{1+\epsilon}}\prod_{p\mid m}\Big(1 - \frac1{p^{1+\epsilon}}\Big)\Big(\sum_r\frac{\mu(r)f(mr)}{r^{1+\epsilon}}\Big)^2,
\]
which follows from $\gcd(d_1, d_2)^{1+\epsilon} = \sum_{m\mid\gcd(d_1,d_2)}m^{1+\epsilon}\prod_{p\mid m}(1 - p^{-1-\epsilon})$.
Since $p^{-1-\epsilon}(1 - p^{-1-\epsilon}) \leq 1/(p - 1)$ and $\zeta(1 + \epsilon) \leq 1/\epsilon + 1$, and $\sum_r|\rho_t(mr)|/r \leq \frac\eta\Delta\beta(m)$, we obtain the bound for $S_t'$; the term $\bar\Phi(\ell_U/\eta)^2$ is not needed but is kept, as in the computation.
For $\Sigma$ we split the sum over $m$.
For $m \leq 10^6$ we have $\ell_K - \log m \geq 26$ for $K \in \{U, V\}$, so $\beta(m) \leq \beta_0$ by Lemma~\ref{lem:A-bsum}, and $\sum_{m\leq y}\mu^2(m)/\varphi(m) = \log y + c_0 + O^*(4/\sqrt y)$ with $c_0 = 1.33258\ldots < 1.3326$ \cite{RamareAkhilesh}.
For $\log m \in [u_j, u_j + \delta_u]$ the same estimate shows that the total mass of $\mu^2(m)/\varphi(m)$ is at most $\delta_u + 8e^{-u_j/2}$.
For $\log m \geq \ell_V + 60\eta$ both arguments of $b$ are at most $-60\eta$, so $\beta(m) \leq 2(\eta + 2)\phi((\log m - \ell_V)/\eta)$, and the mass on intervals of length $1$ is at most $9$; comparing the sum over such intervals with an integral of the decreasing majorant gives the last term.
\end{proof}

\subsection{Proof of Theorem~\ref{thm:G}}

By \eqref{eq:mink}, Proposition~\ref{prop:A-main} and Lemma~\ref{lem:A-St}, the sum in Theorem~\ref{thm:G} is at most $S_s + 2\sqrt{S_sS_t} + S_t$ with $S_s \leq e^{-1/X} + D_4 + \mathrm{REM}$.
It remains to bound $k_1 = H_c(0)$ and $k_0 = H_c'(0) + \gamma_EH_c(0)$.
At $w = 0$ the integrand of $H_c$ is $g(s) = \mathcal A(0; -s, s)e^{\eta^2s^2}/(\zeta(1 - s)\zeta(1 + s)s^4)$, and differentiation with respect to $w$ multiplies it by $D(s) = \frac{\partial}{\partial w}\log\mathcal A + \eta^2s - \frac{\zeta'}{\zeta}(1 + s) - \frac2s$.
On the segment $\Rp s = \kappa_1$, $|\Ip s| \leq 150$, we compute $\frac1{2\pi i}\int g$ and $\frac1{2\pi i}\int gD$ by rigorous numerical integration, with the Euler product truncated at $p \leq 2000$.
The error from the truncation is at most $(e^\tau - 1)\int|g|$, with $\tau$ obtained from $|\mathcal A_p - 1| \leq 5.2p^{-m}$ ($m = 2 - \kappa_1$) and the bound $\pi(t) \leq 1.26t/\log t$ of Rosser and Schoenfeld \cite{RS62}, and similarly for $D$; here $\int|g| \leq (\pi\int|g|^2(1 + t^2)\,dt)^{1/2}$ by the Cauchy--Schwarz inequality.
On the parts $150 < |\Ip s| \leq T_2$ and on the two horizontal pieces of $\mathcal P$ we bound $|g|$ and $|gD|$ in ball arithmetic on pieces of length $0.02$, using the truncated product and the tail bounds $|\mathcal A/\mathcal A_P - 1| \leq e^{\tau'} - 1$ with $\tau' = 4.4P^{1-m'}/(m' - 1)$ and $m' = 1.62$.
This gives $k_1 \in -0.607063 \pm 0.0217$ and $k_0 \in 0.35026 \pm 0.0711$, hence $k_1 \leq -0.5853$ and $k_0 \leq 0.4214$.
Finally, the hypothesis $X > V$ gives $Y_1 > 0$, which is used in Step~4, and $\ell_U \geq 40$ is used in Lemma~\ref{lem:A-St}; the conditions on $\tau$ and $\log X/\log V$ are those under which the computation was carried out.
\qed

\begin{thebibliography}{99}

\bibitem{Axler2} C. Axler, \emph{Estimates for $\pi(x)$ for large values of $x$ and Ramanujan's prime counting inequality}, Integers \textbf{18} (2018), Paper No.~A61.

\bibitem{Axler} C. Axler, \emph{On Ramanujan's prime counting inequality}, preprint (2022), available at \url{https://arxiv.org/abs/2207.02486}.

\bibitem{Bellotti} C. Bellotti, \emph{An explicit log-free zero density estimate for the Riemann zeta-function}, preprint (2024), available at \url{https://arxiv.org/abs/2405.12545}.

\bibitem{BTY} C. Bellotti, T. Trudgian and A. Yang, \emph{Zero-free regions inspired by work of Heath-Brown}, preprint (2026), available at \url{https://arxiv.org/abs/2603.21490}.

\bibitem{BerndtIV} B. C. Berndt, \emph{Ramanujan's Notebooks. Part IV}, Springer-Verlag, New York, 1994.

\bibitem{Buthe2016} J. B\"uthe, \emph{Estimating $\pi(x)$ and related functions under partial RH assumptions}, Math. Comp. \textbf{85} (2016), no.~301, 2483--2498.

\bibitem{Buthe2018} J. B\"uthe, \emph{An analytic method for bounding $\psi(x)$}, Math. Comp. \textbf{87} (2018), no.~312, 1991--2009.

\bibitem{CullyHugillJohnston} M. Cully-Hugill and D. R. Johnston, \emph{On the error term in the explicit formula of Riemann--von Mangoldt}, Int. J. Number Theory \textbf{19} (2023), no.~6, 1205--1228.

\bibitem{De1} S. De, \emph{On proving an inequality of Ramanujan using explicit order estimates for the Mertens function}, preprint (2024), available at \url{https://arxiv.org/abs/2407.12052}.

\bibitem{De2} S. De, \emph{On proving Ramanujan's inequality using a sharper bound for the prime counting function $\pi(x)$}, preprint (2024), available at \url{https://arxiv.org/abs/2408.02591}.

\bibitem{DudekPlatt} A. W. Dudek and D. J. Platt, \emph{On solving a curious inequality of Ramanujan}, Exp. Math. \textbf{24} (2015), no.~3, 289--294.

\bibitem{HSW} E. Hasanalizade, Q. Shen and P.-J. Wong, \emph{Counting zeros of the Riemann zeta function}, J. Number Theory \textbf{235} (2022), 219--241.

\bibitem{Hassani} M. Hassani, \emph{On an inequality of Ramanujan concerning the prime counting function}, Ramanujan J. \textbf{28} (2012), no.~3, 435--442.

\bibitem{Hiary} G. A. Hiary, \emph{An explicit van der Corput estimate for $\zeta(1/2+it)$}, Indag. Math. (N.S.) \textbf{27} (2016), no.~2, 524--533.

\bibitem{HPY} G. A. Hiary, D. Patel and A. Yang, \emph{An improved explicit estimate for $\zeta(1/2+it)$}, J. Number Theory \textbf{256} (2024), 195--217.

\bibitem{arb} F. Johansson, \emph{Arb: efficient arbitrary-precision midpoint-radius interval arithmetic}, IEEE Trans. Comput. \textbf{66} (2017), no.~8, 1281--1292.

\bibitem{Johnston} D. R. Johnston, \emph{Improving bounds on prime counting functions by partial verification of the Riemann hypothesis}, Ramanujan J. \textbf{59} (2022), no.~4, 1307--1321.

\bibitem{JohnstonYang} D. R. Johnston and A. Yang, \emph{Some explicit estimates for the error term in the prime number theorem}, J. Math. Anal. Appl. \textbf{527} (2023), no.~2, Paper No.~127460.

\bibitem{KLN} H. Kadiri, A. Lumley and N. Ng, \emph{Explicit zero density for the Riemann zeta function}, J. Math. Anal. Appl. \textbf{465} (2018), no.~1, 22--46.

\bibitem{MV73} H. L. Montgomery and R. C. Vaughan, \emph{The large sieve}, Mathematika \textbf{20} (1973), 119--134.

\bibitem{Neubauer} G. Neubauer, \emph{Eine empirische Untersuchung zur Mertensschen Funktion}, Numer. Math. \textbf{5} (1963), 1--13.

\bibitem{PY} D. Patel and A. Yang, \emph{An explicit sub-Weyl bound for $\zeta(1/2+it)$}, J. Number Theory (2024).

\bibitem{PlattTrudgian} D. J. Platt and T. S. Trudgian, \emph{The error term in the prime number theorem}, Math. Comp. \textbf{90} (2021), no.~328, 871--881.

\bibitem{PlattTrudgianRH} D. J. Platt and T. S. Trudgian, \emph{The Riemann hypothesis is true up to $3\cdot10^{12}$}, Bull. Lond. Math. Soc. \textbf{53} (2021), no.~3, 792--797.

\bibitem{RamareAkhilesh} O. Ramar\'e and P. Akhilesh, \emph{Explicit averages of non-negative multiplicative functions: going beyond the main term}, Colloq. Math. \textbf{147} (2017), no.~2, 275--313.

\bibitem{RS62} J. B. Rosser and L. Schoenfeld, \emph{Approximate formulas for some functions of prime numbers}, Illinois J. Math. \textbf{6} (1962), 64--94.

\end{thebibliography}

\end{document}
